\documentclass[11pt,a4paper]{article}

% ---------------------------------------------------------------------------
% Copyright (c) 2026 Ashutosh Sinha <ajsinha@gmail.com>. All rights reserved.
% Proprietary and confidential. No licence is granted except by separate
% written agreement. See LICENSE and NOTICE at the root of the repository.
% ---------------------------------------------------------------------------

\usepackage[utf8]{inputenc}
\usepackage[T1]{fontenc}
\usepackage[margin=1.05in]{geometry}
\usepackage{amsmath,amssymb,amsthm}
\usepackage{mathtools}
\usepackage{booktabs}
\usepackage{array}
\usepackage{longtable}
\usepackage{enumitem}
\usepackage{microtype}
\usepackage{xcolor}
\usepackage[most]{tcolorbox}
\usepackage{xurl}
\usepackage{aliascnt}
\usepackage[hidelinks,breaklinks]{hyperref}
\usepackage{cleveref}
\usepackage[htt]{hyphenat}
% Module paths and test names are long unbroken \texttt runs. Let them break after
% an underscore, a slash, a dot or a colon rather than run into the margin.
\let\oldunderscore\_
\renewcommand{\_}{\oldunderscore\allowbreak}
\emergencystretch=3em

\newcolumntype{L}[1]{>{\raggedright\arraybackslash}p{#1}}
% Every environment shares the theorem counter, through an alias, so that cleveref
% names each by its own kind ("Proposition 3.5") rather than by the shared counter.
\newcommand{\sharedtheorem}[3]{%
  \newaliascnt{#1}{theorem}%
  \newtheorem{#1}[#1]{#2}%
  \aliascntresetthe{#1}%
  \crefname{#1}{#2}{#3}\Crefname{#1}{#2}{#3}}
\theoremstyle{plain}
\newtheorem{theorem}{Theorem}[section]
\crefname{theorem}{Theorem}{Theorems}\Crefname{theorem}{Theorem}{Theorems}
\sharedtheorem{proposition}{Proposition}{Propositions}
\sharedtheorem{lemma}{Lemma}{Lemmas}
\sharedtheorem{corollary}{Corollary}{Corollaries}
\theoremstyle{definition}
\sharedtheorem{definition}{Definition}{Definitions}
\sharedtheorem{example}{Example}{Examples}
\sharedtheorem{observation}{Observation}{Observations}
\theoremstyle{remark}
\sharedtheorem{remark}{Remark}{Remarks}

\newcommand{\ie}{\emph{i.e.}\ }
\newcommand{\eg}{\emph{e.g.}\ }
\newcommand{\defn}[1]{\emph{#1}}
\newcommand{\PASS}{\textsf{pass}}
\newcommand{\FAIL}{\textsf{fail}}
\newcommand{\IND}{\textsf{indeterminate}}
\newcommand{\ERR}{\textsf{error}}
\newcommand{\SKIP}{\textsf{skipped}}
\newcommand{\Gam}{\Gamma}
\newcommand{\Dec}{\mathcal{D}}
\newcommand{\Ctl}{\mathcal{C}}
\newcommand{\Ver}{\mathsf{V}}
\newcommand{\hash}{\mathsf{H}}
\newcommand{\pql}[1]{\texttt{#1}}

% ---- expository environments ---------------------------------------------
\definecolor{plainbg}{RGB}{246,246,244}
\definecolor{plainbar}{RGB}{120,120,130}
\definecolor{exbg}{RGB}{247,244,239}
\definecolor{exbar}{RGB}{150,120,70}

\newtcolorbox{plain}{
  enhanced, breakable, sharp corners, boxrule=0pt, frame hidden,
  colback=plainbg, borderline west={2pt}{0pt}{plainbar},
  left=10pt, right=8pt, top=6pt, bottom=6pt,
  before skip=8pt, after skip=8pt,
  fontupper=\small
}
\newtcolorbox{worked}[1][]{
  enhanced, breakable, sharp corners, boxrule=0pt, frame hidden,
  colback=exbg, borderline west={2pt}{0pt}{exbar},
  left=10pt, right=8pt, top=6pt, bottom=6pt,
  before skip=8pt, after skip=8pt,
  fontupper=\small, #1
}
\newcommand{\pt}{\textbf{In plain terms.}\ }
\newcommand{\wk}[1]{\textbf{Worked example: #1}\par\smallskip}

% ---- where a claim is enforced --------------------------------------------
% The argument of this paper is that a quality claim must be justified, and that a
% claim nobody checks is not. A paper is a set of claims, so the same standard
% applies to it: every formal claim below carries a marker saying what carries it
% in the implementation, and the markers that say "nothing" are the ones that
% matter.
\newcommand{\runs}[1]{\par\smallskip\noindent{\footnotesize\raggedright\textsf{Runs:}
  \texttt{#1}\par}\smallskip}
\newcommand{\stated}[1]{\par\smallskip\noindent{\footnotesize\raggedright\textsf{Not
  executed:} #1\par}\smallskip}
\newcommand{\runsprose}[1]{\par\smallskip\noindent{\footnotesize\raggedright\textsf{Runs:}
  #1\par}\smallskip}
\newcommand{\absent}[1]{\par\smallskip\noindent{\footnotesize\raggedright\textsf{Not in
  \textsc{Prama}:} #1\par}\smallskip}
\newcommand{\partly}[1]{\par\smallskip\noindent{\footnotesize\raggedright\textsf{In part:}
  #1\par}\smallskip}

\hypersetup{pdftitle={Data Quality as Justified Belief},
  pdfauthor={Ashutosh Sinha},
  pdfsubject={Derived controls, deterministic verdicts and evidence that verifies without its author},
  pdfkeywords={data quality, semantic layer, rule languages, reconciliation, lineage,
    evidence ledgers, hash chains, conformal prediction, false discovery rate}}

\title{\bfseries Data Quality as Justified Belief\\[6pt]
{\large Derived Controls, Deterministic Verdicts,\\
and Evidence that Verifies Without Its Author}}

\author{Ashutosh Sinha\\[2pt]
{\normalsize\itshape Independent Researcher}\\[2pt]
{\normalsize\ttfamily ajsinha@gmail.com}}
\date{}

\begin{document}

\maketitle

\begin{center}
\footnotesize Copyright (c) 2026 Ashutosh Sinha. All rights reserved. Proprietary and confidential.
\end{center}

\begin{abstract}
\noindent
A statement that a dataset is ``of good quality'' is a belief about data. An organisation acts on the
belief --- it files a regulatory return, closes its books, prices a portfolio --- and is later asked why it
held it. The tools sold to answer that question run on three things that are each typed in by somebody: the
controls (a person writes the SQL), the verdicts (a threshold, and increasingly a model, decides), and the
record (a log that the system producing it also verifies). Each is a place where a belief can be held for
no good reason while looking as though it has one.

We take the classical epistemological requirement --- that knowledge be both true and \emph{justified} --- and
turn it into three engineering constraints. \emph{A control must derive from a declaration}: an owner states,
in business terms, what a dataset's grain is, what values an attribute may take, and how datasets relate,
and a deterministic, total, provenance-carrying map $\Gam$ produces the controls. What no declaration
covers is named by a coverage measure rather than hidden by a percentage of columns touched. \emph{A verdict
must come from a deterministic engine}: controls are written once in a rule language (PQL), lowered to one
engine-neutral intermediate representation, compiled to three SQL engines and checked against an
independent reference interpreter; the verdict set is larger than $\{\PASS,\FAIL\}$ because a SQL screen
that finds nothing has established a lower bound, not a pass; and no code path lets a model's output decide
a verdict. \emph{Evidence must verify without trusting its author}: every judgement is appended to a
hash-chained ledger with a Merkle root, and a verifier written against the standard library alone,
importing nothing from the system, reaches the same conclusion.

On that base we give formal accounts of reconciliation with tolerances and timing offsets, of column
lineage read from code (parsed, never executed) and the controls it can and cannot justify, of trust
propagation as a choice of combining operations (only one of which is, strictly, a semiring), and of
false-alarm budgets as conformal $p$-values with
false-discovery-rate selection. Every formal claim is marked with what carries it in \textsc{Prama}, the
system the author built to this account: a named test, prose, or nothing. The markers are not uniformly
favourable. Two results are negative and are reported in the body: a lineage graph of value flow alone
cannot see a defect that acts through a join predicate (so \textsc{Prama} now records a join as a separate,
row-population edge), and on its own labelled defect benchmark \textsc{Prama}'s declared path finds 10 of 28
planted defects and none in the relational, temporal or semantic families. \Cref{sec:register} counts what runs.
\end{abstract}

\tableofcontents
\newpage

% =============================================================================
\section{Introduction}\label{sec:intro}

\subsection{The domain}

A large bank holds its authoritative data in several hundred places: warehouses, operational databases,
files landed nightly by counterparties, message streams, and the general ledger. Some of that data feeds a
regulatory return whose errors are fined by the million \cite{ateam2024}; some feeds a risk number that a
board relies on; some feeds the month-end close. The Basel Committee's principles for risk data aggregation
\cite{bcbs239} ask a bank to show that such data is accurate, complete and timely, and --- the part that is
expensive --- to \emph{show} it, to a supervisor, after the fact.

The tooling market for this problem is large and fragmented \cite{gartner2025adq,datakitchen2026}. Rule
frameworks let an engineer write assertions against a table \cite{deequ,tfdv}; observability products fit
statistical monitors to column metrics; catalogues record ownership and glossary terms; reconciliation is
sold separately. A research literature on constraint discovery \cite{metanome,fdeval,dcdiscovery} and error
detection and repair \cite{holoclean,rahabaran} addresses what an engineer might check, and a recent one
addresses what a language model might propose \cite{llmdqr,llmtabular,tuwien,zeroed}.

\subsection{Three places a belief is held without reason}

We read the practical failures of this tooling as three failures of justification.

\paragraph{The control is typed in.} A person writes \pql{SELECT count(*) \ldots\ WHERE isin IS NULL}. The
control encodes a belief about what the dataset should be, but the belief is recorded nowhere except in
the control. When the dataset's grain changes, nothing connects the change to the control. When an auditor
asks why a critical data element has no accuracy check, the answer is that nobody wrote one, and nothing
counted the absence.

\paragraph{The verdict is decided by something that is not accountable.} A control ``passes'' because a SQL
screen returned zero rows, when the screen could only express the shape of an identifier and not its check
digit; a freshness control ``passes'' because a row count was positive. Increasingly, a model is asked
whether a column ``looks right''. In each case a verdict is issued by a mechanism that cannot justify it.

\paragraph{The evidence certifies itself.} The execution log lives in the same database, under the same
administrator, as the system whose correctness it records. A record that can be verified only by the
software that wrote it is testimony, not evidence.

\subsection{The thesis}\label{sec:thesis}

The Sanskrit \emph{pram\=a} names valid cognition: knowledge that is both true and arrived at by a reliable
means. A true belief held for a bad reason is not \emph{pram\=a}. We think this is the right standard for a
data quality claim, and that it decomposes into three engineering constraints, each checkable.

\begin{enumerate}[leftmargin=2em,label=\textbf{J\arabic*.}]
  \item \textbf{A control derives from a declaration.} Every control either is produced by a total,
        deterministic map $\Gam$ from something an accountable owner declared, or carries the provenance
        of the proposal that a named person accepted. What is declared but not controlled, and what is
        critical but not declared, is measured.
  \item \textbf{A verdict comes from a deterministic engine.} A verdict is a function of the compiled plan
        and the data, computed the same way on every engine the plan targets, and never a function of a
        model's output. Where the engine cannot decide, the verdict says so.
  \item \textbf{Evidence verifies without its author.} The record of every verdict can be checked by a
        program that shares no code with the system that produced it.
\end{enumerate}

\begin{plain}
\pt ``Declare it. Prove it. Trust it.'' A business owner says what the data means; the system turns that
into checks nobody had to write; a deterministic engine runs them and records what it found in a form that
anyone can check with a short script. A model may help write a check or explain a result. It never decides
whether the data passed.
\end{plain}

\subsection{Contributions}

\begin{enumerate}[leftmargin=2em]
  \item A formal account of \emph{declaration-derived controls}: the map $\Gam$ from a business semantic
        layer (grain, rhythm, value domains, semantic types, and thirteen relationship kinds) to controls,
        with determinism, identity stability, provenance and refusal properties, and a coverage measure over
        (attribute, dimension) pairs (\Cref{sec:gamma}).
  \item A verdict semantics for an engine-neutral rule language that separates an exact test from a
        \emph{screen}, makes $\IND$ first class, and is checked across three SQL engines and a reference
        interpreter; and the fusion of controls into shared scans (\Cref{sec:pql}).
  \item An evidence ledger whose verification is independent of the system that writes it
        (\Cref{sec:evidence}).
  \item A tolerance and timing semantics for reconciliation, with a counting result for timing offsets and
        its measurement on a month-end close (\Cref{sec:recon}).
  \item An account of column lineage read from code, of the controls it justifies and the ones it must
        hold back, a negative result about what value lineage cannot see, and the population edge that
        repairs it (\Cref{sec:lineage}).
  \item The boundary ``models author, the engine decides'', enforced by an architectural test with a
        counterfactual, by a gate that refuses a control that cannot fail, and by an evaluation-gated
        model gateway with a hash-chained call ledger (\Cref{sec:ai}).
  \item Trust propagation as a choice of semiring, and false-alarm budgets as conformal $p$-values with
        Benjamini--Hochberg selection (\Cref{sec:trust}); metadata-implied rules, correlation and fitness
        search (\Cref{sec:metadata}).
  \item An evaluation that reports eight planted-defect case studies, both columns, and a labelled
        benchmark whose bounds and ablations were run and whose comparison with \textsc{Prama} was not
        (\Cref{sec:eval}).
\end{enumerate}

\subsection{How to read this paper, and how to check it}

Each definition, proposition and theorem is followed by a marker:
\textsf{Runs:} names a test in the repository, as \texttt{path::test\_name}, that exercises the claim;
\textsf{Runs:} followed by prose names a mechanism without a single test;
\textsf{In part:} says what is weaker in the code than on the page;
\textsf{Not executed:} marks mathematics that constrains the system without the system checking it; and
\textsf{Not in \textsc{Prama}:} marks a construction that is not built. Every cited test was located in
the repository at the time of writing. A reader with the repository can run any of them with
\texttt{pytest -q '<path>::<name>'}. Numbers are quoted only where they come from a run we performed or from
an assertion in a named test, and each says which.

The system is \textsc{Prama} version 0.1.0 (\texttt{src/prama/version.py}), about 115,000 lines of Python
under \texttt{src/prama} and 249 test modules under \texttt{tests}. It is proprietary; the paper describes
it and does not release it.

% =============================================================================
\section{The objects, in engineering terms}\label{sec:objects}

Before the formal material, the vocabulary in the terms an engineer uses.

\begin{description}[leftmargin=1.5em,style=nextline]
  \item[Estate] Everything a tenant governs: connectors to sources, and the datasets read through them.
  \item[Dataset] A table, a view, a feed of files, or a set of feeds, with an owner, a criticality tier,
        a \emph{grain} (``one row per trade'') and a \emph{rhythm} (``arrives by 06:30 on business days'').
  \item[Attribute] A column given a business meaning: optional or mandatory, a semantic type such as ISIN
        or LEI, a value domain (a code list or a range), and possibly the mark of a \emph{critical data
        element} (CDE).
  \item[Relationship] A typed, declared link between two datasets: \textsf{references},
        \textsf{reconciles-with}, \textsf{derives-from}, and ten others (\Cref{tab:relkinds}).
  \item[Control] An executable check, written in PQL, with a severity, a quality dimension and a
        threshold. A control has a stable \emph{identity} and a content \emph{hash}.
  \item[Proposal] A candidate control, from any origin, awaiting a person. Proposals are the only way a
        control enters the estate.
  \item[Plan] The engine-neutral intermediate representation (IR) a control lowers to; the thing that is
        compiled, fused, executed and hashed.
  \item[Verdict] The outcome of executing a plan: $\PASS$, $\FAIL$, $\IND$, $\ERR$ or $\SKIP$.
  \item[Evidence record] An append-only, hash-linked record of one verdict and what it was computed from.
\end{description}

The pipeline is: \emph{declare} $\to$ $\Gam$ \emph{derives} proposals $\to$ a person \emph{activates}
$\to$ the plan is \emph{compiled} and \emph{fused} $\to$ an engine \emph{executes} $\to$ the verdict is
\emph{judged} from metrics $\to$ the evidence is \emph{recorded} and can be \emph{verified} and
\emph{replayed}. Models sit beside this pipeline --- they may draft a declaration, a proposal or an
explanation --- and never inside the step from metrics to verdict.

% =============================================================================
\section{Controls derived from declarations}\label{sec:gamma}

\subsection{Why derive}

A control written by hand restates a belief about a dataset in a second place. The first place is
wherever the belief really lives: in the owner's head, in a data dictionary, in a regulatory instruction.
Two statements of one rule diverge, and they diverge silently, because only one of them is executed. The
remedy we propose is not better discipline in writing controls but a change in what is written: the owner
writes the belief once, in business terms, and the control is computed from it. When the belief changes,
the control changes, because it is the same object seen twice.

\subsection{Declarations}

\begin{definition}[Dataset declaration]\label{def:decl}
A \defn{dataset declaration} is a tuple
\[ d = (\mathit{ref}, G, K, \rho, \theta, A, \alpha, t) \]
where $\mathit{ref}$ names a dataset; $G$ is a \emph{grain}, a set of attribute names whose values identify a row;
$K$ is an optional business key; $\rho$ is a \emph{rhythm} (an arrival expectation with a calendar, and an
expected volume); $\theta$ is a temporality; $A$ is a finite set of \emph{attribute declarations}; $\alpha$
is an optional authoritativeness statement (``replica of $X$''); and $t$ is a criticality tier. An attribute
declaration $a \in A$ carries an optionality, an optional semantic type (\eg ISIN, LEI, ISO date, monetary
amount with its currency column), an optional value domain (a code list or a range), and a flag marking it
a critical data element. A declaration records who declared it and when.
\end{definition}

\begin{definition}[Relationship declaration]\label{def:rel}
A \defn{relationship declaration} is a tuple
\[ r = (k, \mathit{left}, \mathit{right}, M, c, \epsilon, \delta) \]
where $k$ is one of the thirteen kinds of \Cref{tab:relkinds}, $M$ is a match key (pairs of attributes), $c$ an
optional compared attribute, $\epsilon$ an optional tolerance and $\delta$ an optional timing offset.
A relationship is either \emph{proposed} (discovered from data or lineage) or \emph{confirmed} by a person.
\end{definition}

\runs{tests/semantic/test\_services.py::TestRelationshipService::test\_every\_kind\_is\_offered\_in\_business\_language}
\runs{tests/semantic/test\_services.py::TestRelationshipService::test\_a\_discovered\_relationship\_is\_never\_confirmed\_automatically}

\subsection{The map $\Gam$}

\begin{definition}[Control, identity and hash]\label{def:control}
A \defn{control} is a PQL statement together with a severity, a quality dimension, and a
\emph{provenance}. Its \defn{identity} is a function $\iota(\mathit{declaration}, \mathit{rule},
\mathit{subject})$ of where it came from and never of its text. Its \defn{content hash} is a function of
its text. Identity answers ``is this the same control?''; the hash answers ``has it changed?''.
\end{definition}

\begin{definition}[$\Gam$]\label{def:gamma}
For a dataset declaration $d$, $\Gam(d)$ is obtained by applying, in a fixed order, the rules
$\gamma_{\mathrm{grain}}$, $\gamma_{\mathrm{key}}$, $\gamma_{\mathrm{rhythm}}$, $\gamma_{\mathrm{temporality}}$,
$\gamma_{\mathrm{attr}}$ and $\gamma_{\mathrm{auth}}$, each of which returns a finite set of controls and a
finite set of \emph{unsatisfiable} findings, and then \emph{coalescing}: controls whose checks are equal are
merged into one control whose rule is the sorted join of the contributing rules. For a relationship $r$,
$\Gam(r)$ dispatches on the kind $k$ and returns a lineage edge together with either controls,
comparison specifications for the reconciliation engine, or a refusal naming what is missing.
\end{definition}

\Cref{tab:gamma} lists what each dataset rule produces and \Cref{tab:relkinds} what each relationship
kind produces. The second table is honest about a gap between what a kind \emph{declares} it generates
and what the code \emph{emits}: every kind emits something, but several emit fewer control families than
the design corpus lists.

\begin{table}[t]
\centering\small
\begin{tabular}{@{}L{0.30\linewidth}L{0.62\linewidth}@{}}
\toprule
\textbf{Declaration} & \textbf{Control emitted by $\Gam$} \\
\midrule
Grain $G$ & \pql{HAS UNIQUE KEY} over $G$ (\textsf{grain.uniqueness}); completeness of each member of $G$ \\
Business key $K$ & uniqueness over $K$ \\
Rhythm $\rho$ & a freshness control with the declared calendar; a volume control; seasonality deferred to a monitor \\
Temporality $\theta$ & a temporal control \\
Mandatory attribute & \pql{IS NOT NULL} \\
Semantic type & \pql{IS VALID '<type>'}, naming its standard \\
Value domain & \pql{IN CODELIST} or a range \\
Monetary amount & \pql{IN CODELIST 'iso4217'} on the currency column beside it \\
Authoritativeness & parity with the authoritative source \\
\bottomrule
\end{tabular}
\caption{The dataset rules of $\Gam$ (\texttt{src/prama/derive/generator.py}).}\label{tab:gamma}
\end{table}

\begin{table}[t]
\centering\small
\begin{tabular}{@{}L{0.22\linewidth}L{0.34\linewidth}L{0.36\linewidth}@{}}
\toprule
\textbf{Kind} & \textbf{Declared families} & \textbf{Actually emitted} \\
\midrule
references & referential integrity, orphans, key coverage & PQL \pql{REFERENCES} (a correlated \pql{EXISTS}) \\
reconciles-with & reconciliation, breaks, certificate & a reconciliation specification \\
derives-from & aggregate parity, trust edge, impact & aggregate parity \\
feeds & lineage edge, arrival chain, latency & an edge, and deliberately no control \\
mirrors & row-count, content parity, staleness & row-count parity, value parity (if compared), staleness \\
aggregates & roll-up parity, trust edge & aggregate parity \\
enriches & enrichment coverage & PQL reference check, completeness dimension \\
supersedes & migration parity, dual run & value parity \\
same-entity-as & resolution, duplicates, identifiers & identifier consistency \\
temporal-successor & roll-forward & roll-forward \\
parent-of & hierarchy completeness, cycles, orphans & PQL orphan-node check only; no cycle detection \\
mutually-exclusive & overlap & overlap \\
together-complete & population completeness & coverage \\
\bottomrule
\end{tabular}
\caption{The thirteen relationship kinds (\texttt{src/prama/derive/relationships.py}). Only three kinds
yield PQL controls; the rest yield comparison specifications executed by the reconciliation engine.}
\label{tab:relkinds}
\end{table}

\begin{proposition}[Determinism and identity stability]\label{prop:det}
$\Gam$ is a function: regenerating an unchanged declaration yields the same identities and the same
hashes. Editing a declaration so that a control's check changes yields a control with the same identity
and a different hash, so the control is updated rather than orphaned and replaced.
\end{proposition}
\begin{proof}
Each rule is a pure function of the declaration, applied in a fixed order, and coalescing sorts its rule
names, so the output is independent of iteration order. Identity is computed from $(\mathit{declaration},
\mathit{rule}, \mathit{subject})$, none of which an edit to a threshold or a code list changes; the hash is
computed from the text, which such an edit does change.
\end{proof}
\runs{tests/derive/test\_generator.py::test\_regenerating\_an\_unchanged\_declaration\_changes\_nothing}
\runs{tests/derive/test\_generator.py::test\_editing\_a\_declaration\_updates\_a\_control\_rather\_than\_orphaning\_it}
\runs{tests/propose/test\_end\_to\_end.py::test\_a\_second\_run\_of\_the\_generator\_adds\_nothing}

\begin{proposition}[Provenance]\label{prop:prov}
Every control in $\Gam(d)$ carries a provenance naming its origin (\textsf{declaration}), the rule that
produced it, the declaration and attribute it came from, the declaring person and the time. The question
``why does this control exist?'' therefore has an answer that is data, not recollection.
\end{proposition}
\runs{tests/derive/test\_generator.py::test\_every\_generated\_control\_can\_say\_why\_it\_exists}

\begin{proposition}[Round trip]\label{prop:roundtrip}
Every generated control is valid PQL: for every $c \in \Gam(d)$, $\mathrm{render}(\mathrm{parse}(c)) = c$.
\end{proposition}
\runs{tests/derive/test\_generator.py::test\_every\_generated\_control\_re\_parses\_to\_itself}

\begin{proposition}[Coalescing]\label{prop:coalesce}
If two rules produce the same check, $\Gam$ emits one control, whose rule records both.
\end{proposition}
\runs{tests/derive/test\_generator.py::test\_two\_rules\_reaching\_the\_same\_check\_produce\_one\_control}

\begin{proposition}[Refusal, not omission]\label{prop:refusal}
$\Gam$ is total in the following sense: every declaration either produces a control or produces a finding
that says why it did not. A grain naming a column that does not exist is reported, not emitted; a
reconciliation with no tolerance is refused with the reason; every one of the thirteen relationship kinds
yields an edge and either an output or a refusal.
\end{proposition}
\runs{tests/derive/test\_generator.py::test\_a\_grain\_naming\_a\_column\_that\_does\_not\_exist\_is\_reported\_not\_emitted}
\runs{tests/derive/test\_relationships.py::test\_a\_reconciliation\_with\_no\_tolerance\_is\_refused\_with\_the\_reason}
\runs{tests/derive/test\_relationships.py::test\_every\_kind\_produces\_something\_and\_nothing\_silently}
\runs{tests/derive/test\_relationships.py::test\_feeds\_generates\_an\_edge\_and\_deliberately\_no\_control}

\begin{plain}
\pt Writing ``one row per trade, identified by trade id'' produces a uniqueness check and a not-null check
on the trade id. Writing ``this column is an ISIN'' produces a validity check that names ISO 6166. Nobody
writes SQL. If the owner changes the rule, the same control is updated. If a statement cannot be turned into
a check --- a grain naming a column that is not there --- the system says so instead of quietly dropping it.
\end{plain}

\begin{worked}
\wk{a trading book}
Case study 1 (\texttt{case-studies/01-trading-book-sqlite}) declares six datasets over one SQLite database,
with 5,000 trades and nine planted defects. In our run, $\Gam$ produced 59 controls and reported nine
declarations that produced none, each with a reason: six declare when a dataset arrives but no column that
records arrival, so there is nothing to measure freshness on, and three name a validator Prama does not have
(\emph{``Prama has no validator or code list called `currency'\,''}). Executed, the 59 controls gave 44
passes, 5 failures and 10 not established. The five
failures are exactly the five planted defects that a single-dataset declaration can see: the duplicated
positions (grain), 60 missing notionals (a mandatory CDE), 7 negative quantities (a range), and the currency
\texttt{'EURO'} (a code list, reported twice, as validity and consistency). The four it did not find are
each named in the run's output with the reason: an ISIN checksum (a screen, \Cref{sec:screen}), an orphan
instrument and a ledger adjustment (both need a \emph{relationship}, which this study does not declare), and
a settlement date before the trade date --- for which no control exists because nothing in the declaration
says how the two dates relate. That last one is a gap in the declaration, not in $\Gam$, and the study
leaves it in to show the difference.
\end{worked}

\subsection{Coverage, and the number that flatters}\label{sec:coverage}

A percentage of columns ``with a control'' flatters: a single \pql{IS NOT NULL} makes a column covered
while its format, domain and accuracy go unchecked.

\begin{definition}[Applicable pairs and coverage]\label{def:coverage}
For a declared dataset, the \defn{applicable pairs} $P$ are the (attribute, dimension) pairs, over the
dimensions completeness, validity, consistency and accuracy, that the declaration makes meaningful:
validity only where a semantic type or a value domain is declared; consistency only for an amount with a
currency; accuracy only for a CDE; completeness except for an optional attribute. Uniqueness and timeliness
are counted at dataset level. The \defn{coverage} is $\kappa = |P_{\mathrm{cov}}|/|P|$, where
$P_{\mathrm{cov}}$ are the pairs some active control addresses. The \defn{touched fraction} $\tau$ is the
share of attributes with at least one control.
\end{definition}

\begin{proposition}[The honest and the flattering number]\label{prop:flatter}
$\tau$ can equal $1$ while $\kappa < 1/2$. Both are reported, side by side, and every uncovered pair is
reported with the declaration or control that would close it, ranked by a risk weight that favours CDEs
and obligations.
\end{proposition}
\begin{proof}
Take attributes each with a mandatory flag, a semantic type and a CDE mark, and one \pql{IS NOT NULL} control
on each. Then $\tau = 1$, while of the three applicable pairs per attribute only completeness is covered,
so $\kappa = 1/3$.
\end{proof}
\runs{tests/derive/test\_coverage.py::test\_coverage\_is\_measured\_per\_dimension\_not\_per\_column}
\runs{tests/derive/test\_coverage.py::test\_the\_flattering\_number\_is\_reported\_beside\_the\_honest\_one}
\runs{tests/derive/test\_coverage.py::test\_a\_dimension\_an\_attribute\_cannot\_have\_is\_not\_counted\_against\_it}
\runs{tests/derive/test\_coverage.py::test\_every\_gap\_names\_what\_would\_close\_it}

\begin{proposition}[A dataset cannot vouch for its own accuracy]\label{prop:accuracy}
No dataset declaration alone covers the accuracy pair of a CDE. Coverage of a CDE's accuracy requires a
relationship declaration.
\end{proposition}
\begin{proof}
The accuracy dimension is satisfied only by a control comparing the attribute with something outside its
row --- in \textsc{Prama}, a \pql{REFERENCES} control or a comparison. The dataset rules of
\Cref{tab:gamma} emit neither; only the relationship rules of \Cref{tab:relkinds} do.
\end{proof}
\runs{tests/propose/test\_wave6\_acceptance.py::test\_a\_dataset\_declaration\_alone\_cannot\_cover\_its\_own\_cdes}
\runs{tests/propose/test\_wave6\_acceptance.py::test\_adding\_the\_relationships\_reaches\_the\_target}

\begin{plain}
\pt Accuracy means ``agrees with something else''. A dataset, however carefully described, cannot be
compared with nothing. That is a small piece of logic, but it is the reason relationships --- a line drawn
between two datasets --- carry the most weight in \textsc{Prama}'s estate-maturity score.
\end{plain}

The implementation roadmap reports the effect on a realistic Tier-1 declaration: the dataset declaration
alone reaches 82\% of pairs and 71\% of CDE pairs, and adding the relationship declarations reaches 100\%
of both. The test asserts the weaker statement that relationships reach at least 95\% of CDE pairs and 80\%
overall; the 82/71 figures are prose and we did not reproduce them.
\partly{the proposition runs; the specific percentages are stated in \texttt{docs/19}, not asserted.}

\subsection{Proposals, and who activates}\label{sec:proposals}

\begin{definition}[Proposal and origin]\label{def:proposal}
A \defn{proposal} is a candidate control with an \emph{origin} in the totally ordered set
$\textsf{declaration} > \textsf{document} > \textsf{import} > \textsf{mining} > \textsf{example} >
\textsf{induction}$ (the last is a language model), a backtest against current data, and a utility score.
Proposals are the only channel through which a control enters the estate.
\end{definition}

\begin{proposition}[Only a declaration may skip review, and not always]\label{prop:review}
A proposal may be marked as not needing review only if its origin is \textsf{declaration}, and not even then
if its backtest violates more than 20\% of rows (a declaration that would ``drown the queue'' is probably
wrong about the data). Mining, documents, examples and model induction always need review.
\end{proposition}
\runs{tests/derive/test\_provenance.py::test\_only\_a\_declaration\_may\_run\_without\_review}
\runs{tests/propose/test\_queue.py::test\_not\_even\_a\_declaration\_may\_run\_if\_it\_would\_drown\_the\_queue}
\runs{tests/induce/test\_inducer.py::test\_a\_good\_answer\_becomes\_an\_induced\_control}

\begin{remark}
Reading the code closely, the ``may skip review'' mark is counted in the queue summary and used nowhere
else: the only activation path is \texttt{ControlDAO.activate}, which requires an \texttt{approved\_by}
identity and writes it into the control's provenance. In practice, therefore, a person activates every
control, including a declared one; the case studies do so explicitly. The generator's docstring (``nothing it
emits is active'') is correct; the \texttt{may\_auto\_activate} property overstates what happens.
\end{remark}
\runsprose{\texttt{src/prama/db/dao/control.py} (\texttt{activate} requires \texttt{approved\_by}); no test
asserts the absence of another activation path.}

\begin{proposition}[Corroboration merges]\label{prop:corrob}
When two origins propose a control with the same identity, the result is one proposal, held by the
stronger origin and marked corroborated, whose confidence rises by a bounded amount. A rejected proposal is
not offered again unless the evidence changes materially.
\end{proposition}
\runs{tests/propose/test\_queue.py::test\_two\_origins\_reaching\_the\_same\_rule\_merge\_rather\_than\_duplicate}
\runs{tests/propose/test\_queue.py::test\_the\_stronger\_origin\_keeps\_the\_proposal}
\runs{tests/propose/test\_wave6\_acceptance.py::test\_mining\_and\_declaring\_the\_same\_rule\_produces\_one\_corroborated\_proposal}
\runs{tests/propose/test\_utility.py::test\_corroboration\_raises\_confidence\_modestly}
\runs{tests/propose/test\_queue.py::test\_a\_rejected\_proposal\_is\_not\_offered\_again}

The utility score is a fixed linear combination of criticality, coverage gap, confidence and an
acceptance prior learned from past decisions, Laplace-smoothed and capped so that the learned term cannot
dominate. Tier-1 datasets require a second person (maker--checker).
\runs{tests/propose/test\_utility.py::test\_the\_learned\_term\_cannot\_dominate\_the\_ranking}
\runs{tests/semantic/test\_services.py::TestApprovalPolicy::test\_tier\_one\_requires\_a\_second\_person}

\subsection{What $\Gam$ does not yet do}

\paragraph{Freshness, and the column it needs.} A declared rhythm says when data is due; it does not say how
anybody would know it arrived, and a plain table keeps its rows, not when they were loaded. Until recently
$\Gam$ generated a freshness control anyway, which had no execution strategy and could never reach a verdict;
the repository pinned it as a \emph{strict} expected failure. It is closed by making the missing fact
declarable. A rhythm names its \emph{arrival column} (a load timestamp), and $\Gam$ derives
\pql{CHECK t.loaded\_at IS FRESH WITHIN 30 MINUTES OF '06:30' CALENDAR 'TARGET2'}. It measures
$\max(\textit{loaded\_at})$. The cycle judged is the most recent business day $d$ whose deadline (due time plus
tolerance, on the calendar) has passed at the instant of evaluation $t$. The control passes iff the newest arrival
is later than the deadline of the business day before $d$. The instant $t$ is recorded as a metric, so a replay
of the evidence reaches the same verdict. A rhythm \emph{without} an arrival column generates no control and an
unsatisfiable entry that says what to declare, which is \Cref{prop:refusal} applied to time. The measure is
one-sided by construction: the newest arrival alone cannot tell a very late delivery for the previous cycle
from an early one for this cycle, and it cannot miss the failure that matters, nothing new since the last
deadline.
\runs{tests/derive/test\_generator.py::test\_a\_generated\_freshness\_control\_can\_reach\_a\_verdict}
\runs{tests/backend/test\_freshness.py::test\_the\_cycle\_due\_is\_found\_on\_the\_business\_calendar}
\runs{tests/backend/test\_freshness.py::test\_the\_verdict\_is\_reproducible\_from\_the\_evidence}
\runs{tests/backend/test\_freshness.py::test\_executed\_on\_a\_real\_engine}

The design corpus lists roll-forward, aggregate parity and row-count parity as \emph{PQL} forms. They exist
as comparison specifications produced by $\Gam$; the PQL surface for them is still a specification.
\absent{PQL syntax for roll-forward, aggregate and row-count parity; cycle detection for
\textsf{parent-of}.}

% =============================================================================
\section{One language, one plan, a verdict that can say ``not known''}\label{sec:pql}

\subsection{PQL and the plan}

PQL is \textsc{Prama}'s control language. A control is a \pql{CHECK} over a dataset --- a predicate over rows
(\pql{amount >= 0}, \pql{ccy IN CODELIST 'iso4217'}, \pql{isin IS VALID 'isin'}), a dataset property
(\pql{HAS UNIQUE KEY (trade\_id)}, \pql{REFERENCES}), a named Python delegate, or a single read-only custom
query --- with a threshold (\pql{AT MOST 5 ROWS}, \pql{BELOW 0.1\%}), a severity, a dimension, an optional
scope (\pql{WHERE}, \pql{FOR EACH}) and a policy for unknowns. A \pql{RECONCILE} statement compares two
datasets (\Cref{sec:recon}).

\begin{verbatim}
CHECK trades.notional IS NOT NULL SEVERITY critical DIMENSION completeness
CHECK trades.isin IS VALID 'isin'
CHECK trades HAS UNIQUE KEY (trade_id)
CHECK trades.ccy IN CODELIST 'iso4217'
      TREAT UNKNOWN AS PASS BECAUSE 'optional on internal moves'
\end{verbatim}

\begin{definition}[Plan and plan identity]\label{def:plan}
A control \emph{lowers} to a \defn{plan}: a typed, engine-neutral logical form comprising an assertion
kind, a predicate, an unknown policy, a scope, the metrics to compute and a threshold, together with any
code-list values and validator implementations it names. Its \defn{meaning} $\mu(P)$ is that tuple and its
IR version; its identity is $\mathit{plan\_id}(P) = \texttt{ir:sha256:}\,\hash(\mathrm{canon}(\mu(P)))$,
where $\mathrm{canon}$ serialises with sorted keys, no whitespace and UTF-8. Severity, description,
rationale and provenance are not part of $\mu$.
\end{definition}

\begin{proposition}[What the plan identity closes over]\label{prop:closure}
Changing a code list changes the plan identity of every control that names it; a plan that needs a
residual validator hashes differently from the same plan without one; and the identity is the same
whichever entry point lowered the control.
\end{proposition}
\runs{tests/backend/test\_worked\_example.py::TestTheCodelistIsClosedOver::test\_changing\_the\_codelist\_changes\_the\_control}
\runs{tests/backend/test\_two\_stage.py::test\_a\_two\_stage\_plan\_hashes\_differently\_from\_a\_screen\_only\_one}
\runs{tests/backend/test\_resolve.py::TestItDoesNotChangeAnythingElse::test\_a\_control\_with\_no\_codelist\_hashes\_identically}
\partly{no test asserts that severity or provenance are \emph{excluded} from the hash; it is true of
\texttt{ControlPlan.meaning} by inspection.}

\begin{plain}
\pt A control's fingerprint changes when what it checks changes --- including the list of allowed values
it points at, and including the Python behind a validator it names --- and does not change when somebody
edits its description. Last month's evidence therefore refers to exactly the control that produced it.
\end{plain}

\subsection{The verdict set}

\begin{definition}[Verdicts]\label{def:verdict}
$\Ver = \{\PASS, \FAIL, \IND, \ERR, \SKIP\}$. Given the metrics $m$ an engine returned for a plan with
threshold $\theta$:
\[
\mathrm{judge}(m,\theta) =
\begin{cases}
\IND & \text{if the metric $\theta$ reads is absent,}\\
     & \text{or $m.\mathit{scanned\_rows}$ is absent or $0$,}\\
     & \text{or $\theta$ is relative and its denominator is $0$;}\\
\PASS & \text{if $\theta$ holds of $m$;}\\
\FAIL & \text{otherwise.}
\end{cases}
\]
$\ERR$ records an execution failure (an unreadable source, an absent delegate host) as a finding rather
than a crash. A segmented control fails if any segment fails, and is indeterminate if any segment is
indeterminate and none fails. The actionable verdicts are $\{\FAIL, \ERR, \IND\}$.
\end{definition}

\begin{proposition}[Nothing scanned is not a pass]\label{prop:empty}
A control over an empty scope is $\IND$, for absolute and relative thresholds alike. The exception is a
row-count control, which is a claim about the count itself: over an empty table it fails.
\end{proposition}
\runs{tests/backend/test\_reference.py::TestScopes::test\_an\_empty\_scope\_is\_indeterminate\_not\_a\_pass}
\runs{tests/backend/test\_two\_stage.py::TestNothingScannedIsNotAPass::test\_an\_absolute\_threshold\_refuses\_an\_empty\_scan}
\runs{tests/backend/test\_two\_stage.py::TestNothingScannedIsNotAPass::test\_the\_two\_agree}
\runs{tests/backend/test\_two\_stage.py::TestNothingScannedIsNotAPass::test\_a\_row\_count\_control\_still\_judges\_an\_empty\_table}
\runs{tests/delegates/test\_delegates.py::test\_without\_a\_delegate\_host\_the\_control\_is\_an\_error\_not\_a\_skip}

\subsection{Unknowns}

SQL evaluates a comparison with a null to \textsf{unknown}, and a \pql{WHERE} clause drops unknown rows. A
control written as ``count the rows where the predicate is false'' therefore silently passes every null.

\begin{definition}[Unknown policy]\label{def:unknown}
For a row predicate $\varphi$ with three-valued semantics, the violation count under policy $\pi$ is
\[
v_\pi = \bigl|\{ x : \varphi(x) = \textsf{false} \}\bigr| + [\pi = \textsf{violation}]\cdot
\bigl|\{ x : \varphi(x) = \textsf{unknown} \}\bigr|.
\]
The default is $\pi = \textsf{violation}$. \pql{TREAT UNKNOWN AS PASS} selects the other policy and must
carry a \pql{BECAUSE} clause stating why.
\end{definition}

\begin{proposition}[The default is the conservative one]\label{prop:unknown}
$v_{\textsf{pass}} \le v_{\textsf{violation}}$, with equality exactly when no row evaluates to unknown.
Hence for any upper-bound threshold, a control that passes under the default also passes under
\pql{TREAT UNKNOWN AS PASS}, and not conversely.
\end{proposition}
\begin{proof}Immediate from \Cref{def:unknown}; the second term is non-negative. An upper-bound threshold is
monotone in $v$.\end{proof}
\runs{tests/pql/test\_parser.py::TestUnknownsAreViolationsByDefault::test\_an\_unknown\_counts\_against\_the\_control}
\runs{tests/pql/test\_parser.py::TestUnknownsAreViolationsByDefault::test\_but\_only\_with\_a\_reason}
\runs{tests/backend/test\_engine\_conformance.py::TestTheCorpusDiscriminates::test\_the\_unknown\_policy\_changes\_the\_count}
\runs{tests/backend/test\_reference.py::TestUnknownsAreViolations::test\_a\_null\_counts\_against\_the\_control}

\subsection{A screen is not a test}\label{sec:screen}

Some validity rules have no faithful SQL form. An LEI's check digits are computed by ISO 7064 MOD 97-10
over a letter-to-number mapping; an ISIN's by a Luhn variant. SQL can check the \emph{shape} (length,
alphabet, position of digits) but not the arithmetic, portably. The tempting move is to compile the shape
check and call it the control. That control passes every well-shaped fabricated identifier.

\begin{definition}[Screen and residual]\label{def:screen}
For an exact validator $V$ over values, a \defn{screen} is a predicate $S$ with $V(x) \Rightarrow S(x)$ for
every $x$: a necessary condition every valid value satisfies, and which invalid values may also satisfy.
The \defn{residual} is $V$ itself, run over the rows the screen admits. A plan whose SQL computes only
$S$ is \emph{incomplete}.
\end{definition}

\begin{proposition}[Screen soundness]\label{prop:screen}
Let $n_S = |\{x : \neg S(x)\}|$ and $n_V = |\{x : \neg V(x)\}|$ over a scope. Then $n_S \le n_V$. Hence an
incomplete plan may soundly report $\FAIL$ when $n_S$ breaches the threshold, and may not report $\PASS$
when $n_S = 0$: the only sound verdict for an incomplete plan with $n_S$ within the threshold is $\IND$.
\end{proposition}
\begin{proof}
By contraposition of $V \Rightarrow S$, $\neg S(x) \Rightarrow \neg V(x)$, so the screen's violations are a
subset of the validator's. A breach by the subset is a breach by the whole, for any upper-bound threshold.
Conversely $n_S = 0$ is consistent with any $n_V \ge 0$, so it cannot establish a threshold on $n_V$.
\end{proof}
\runs{tests/backend/test\_two\_stage.py::test\_a\_valid\_value\_is\_never\_rejected\_by\_the\_screen}
\runs{tests/backend/test\_two\_stage.py::test\_the\_screen\_alone\_would\_have\_passed\_the\_fabricated\_value}
\runs{tests/backend/test\_two\_stage.py::test\_the\_compiled\_query\_declares\_that\_it\_is\_not\_the\_whole\_test}
\runs{tests/execute/test\_control\_run.py::TestAnIncompleteScreenCannotPass::test\_zero\_violations\_from\_a\_screen\_is\_not\_a\_pass}
\runs{tests/execute/test\_control\_run.py::TestAnIncompleteScreenCannotPass::test\_a\_screen\_that\_found\_violations\_still\_fails}

\begin{worked}
\wk{a fabricated LEI}
A twenty-character value with the right alphabet and wrong check digits passes the screen: the SQL counts
$n_S = 0$. The exact validator counts $n_V = 1$. \textsc{Prama} compiles the screen, marks the query
incomplete, and on $n_S = 0$ reports \emph{``a pass cannot be reported from a screen alone''}. Across the
eight case studies this is the most frequent ``not established'' line: in case study 1, ten controls on
ISIN, LEI and ISO-date columns are reported this way, including the one on the column where twelve ISINs
with wrong checksums were planted.
\end{worked}

\begin{plain}
\pt Zero rows failing a rough check is not the same as zero rows failing the real check. Reporting it as a
pass would be the most expensive lie a data quality tool can tell, because it is told about exactly the
columns whose correctness SQL cannot express.
\end{plain}

The residual is not run on the SQL path in the case studies: the run reports ``residual not run''. This is
by design --- the residual is a Python validator and runs where rows are available --- but it means that on
a pure pushdown deployment the checksum columns stay $\IND$ until a residual executor is configured.
\partly{screens are sound and never report a pass; the residual executes in the reference interpreter,
not as part of an SQL pushdown run.}

\subsection{One plan, several engines}\label{sec:conformance}

\textsc{Prama} compiles plans to three SQL dialects --- PostgreSQL, DuckDB and SQLite ---
and executes them with a fourth, independent evaluator: a reference interpreter in Python that is never
given SQL and shares no code with the compiler.

\begin{definition}[Conformance]\label{def:conformance}
A backend $e$ \defn{conforms} on a corpus case $c$ if it returns the verdict, the metrics and the
per-segment breakdown that the case declares, or if it refuses $c$ at compile time with a named reason.
A refusal is a conforming outcome; a silently different answer is not.
\end{definition}

\begin{proposition}[Cross-engine agreement on the corpus]\label{prop:conform}
Every case of the conformance corpus yields the same verdict and metrics on every engine that runs it, and
each case is compared by at least two genuinely different engines. The corpus discriminates: it contains
passing and failing cases, a case where the unknown policy changes the count, a duplicate key that counts
once, a segment that fails alone, and rates judged against the rows actually scanned after a filter.
\end{proposition}
\runs{tests/backend/test\_engine\_conformance.py::TestTheEnginesAgree::test\_every\_case\_gives\_the\_same\_answer\_everywhere}
\runs{tests/backend/test\_engine\_conformance.py::TestTheEnginesAgree::test\_at\_least\_two\_genuinely\_different\_engines\_took\_part}
\runs{tests/backend/test\_engine\_conformance.py::TestTheCorpusDiscriminates::test\_a\_rate\_is\_judged\_against\_the\_rows\_actually\_scanned}
\runs{tests/backend/test\_reference.py::TestItSharesNothingWithTheCompiler::test\_it\_never\_asks\_for\_sql}

The corpus has 25 cases (\texttt{src/prama/backend/corpus.py}). DuckDB, SQLite and the reference
interpreter always run; PostgreSQL runs when \texttt{PRAMA\_TEST\_POSTGRES\_DSN} is set. We ran the backend
suite both ways: without PostgreSQL, and against a PostgreSQL 16 container, where all 138 backend tests
passed with none skipped.

Beyond the corpus, a seeded generator produces random controls (250 by default) and requires every engine to
agree on each; the generated set is required to exercise all three of $\PASS$, $\FAIL$ and $\IND$, and every
generated control is reproducible from its seed.
\runs{tests/backend/test\_generated\_equivalence.py::TestTheEnginesAgreeOnAllOfThem::test\_no\_generated\_control\_produces\_a\_disagreement}
\runs{tests/backend/test\_generated\_equivalence.py::TestTheEnginesAgreeOnAllOfThem::test\_the\_generated\_set\_exercises\_all\_three\_outcomes}
\runs{tests/backend/test\_generated\_equivalence.py::TestEveryGeneratedControlIsValid::test\_they\_are\_reproducible\_from\_their\_seed}

\begin{remark}
This is seeded differential testing, not property-based testing with shrinking, and it is not a proof of
semantic equivalence. Full equivalence across engines is probably unattainable for regular-expression
flavours, collations and decimal semantics; the response is a portable subset and compile-time refusal.
\end{remark}
\stated{a semantic-equivalence proof between the compiler and the reference interpreter.}

\begin{proposition}[Refusal is conformance]\label{prop:refuse}
A construct an engine cannot express is refused at compile time with a named reason, never approximated.
SQLite refuses approximate distinct counts and sampling, and a key wider than its column limit; it does
regular expressions through a function the executor registers, and the corpus confirms it agrees with
DuckDB. A fused group containing a refused control is stopped, not silently shortened.
\end{proposition}
\runs{tests/backend/test\_engine\_conformance.py::TestRefusalIsConformance::test\_sqlite\_refuses\_what\_it\_genuinely\_cannot\_do}
\runs{tests/backend/test\_engine\_conformance.py::TestRefusalIsConformance::test\_sqlite\_does\_regex\_through\_a\_registered\_function}
\runs{tests/pql/test\_deep\_input.py::test\_a\_500\_column\_key\_runs\_where\_the\_engine\_can\_and\_is\_refused\_where\_it\_cannot}
\runs{tests/backend/test\_fuse.py::TestTheFusedQuery::test\_a\_control\_the\_engine\_cannot\_run\_stops\_that\_group}
\partly{DuckDB's refusal of back-references and look-arounds (RE2) is implemented and untested.}

\begin{plain}
\pt The same control means the same thing on every database it runs on, and a fourth, deliberately
separate implementation checks that. Where a database genuinely cannot do something, the control refuses
to run there and says why; it never runs a weaker version and reports the answer as if it were the same.
\end{plain}

Building the case studies against real data found two defects of exactly this kind, both since fixed and
recorded in \texttt{case-studies/README.md}: a rounding function that double-rounded on DuckDB (a bare
\pql{CAST(x AS NUMERIC)} is \pql{DECIMAL(18,3)} there), producing 131 false one-cent breaks in 3,000 rows; and
an unknown function name that compiled straight through to SQL while the reference interpreter returned
unknown. The second is the more instructive: the independent check existed, and the two disagreed silently
because nothing compared them on that input. There is now a function catalogue in which no function can
exist without both a lowering and a reference implementation.

\subsection{Fusion}

A source pays for scans, not for controls.

\begin{definition}[Fusion]\label{def:fusion}
Plans over the same dataset and scope are grouped into one query that computes every plan's metrics in a
single pass, computing identical metrics once. The cost of a suite is the number of scans.
\end{definition}

\begin{proposition}[Fusion preserves answers]\label{prop:fusion}
Every control in a fused query receives the metrics, and hence the verdict, it receives when run alone.
\end{proposition}
\begin{proof}[Proof sketch]
Each plan's metrics are aggregate expressions over the same scope, conditioned on the plan's own predicate
(\eg a \pql{SUM(CASE WHEN \ldots)}); aggregates over one scan are independent of the other aggregates
computed in it. Plans whose scopes differ are not fused.
\end{proof}
\runs{tests/backend/test\_fuse.py::TestTheFusedQuery::test\_it\_answers\_every\_control\_identically}
\runs{tests/soak/test\_throughput.py::TestFusionActuallySaves::test\_fused\_answers\_match\_unfused\_ones\_at\_volume}
\runs{tests/backend/test\_fuse.py::TestCost::test\_cost\_is\_counted\_in\_scans\_not\_controls}

The unit suite runs nine controls in three scans. At volume, the soak test asserts that a realistic
generated suite needs at most half the naive number of scans and that the saving grows with the suite; in our
run it printed \emph{``400 controls in 160 scans = 40.0\% of a naive pass''}.
\runs{tests/soak/test\_throughput.py::TestFusionActuallySaves::test\_a\_realistic\_suite\_needs\_far\_fewer\_scans}
\runs{tests/soak/test\_throughput.py::TestFusionActuallySaves::test\_the\_saving\_grows\_with\_the\_suite}

\subsection{Where PQL stops: custom SQL and delegates}

Two escape hatches exist, and both are shaped to keep the verdict with the engine.

\pql{CHECK CUSTOM SQL} accepts exactly one read-only query, checked at parse time, and runs only on the
engines it names. There is deliberately no inline Python.
\runs{tests/pql/test\_custom\_sql.py::test\_anything\_but\_one\_read\_only\_query\_is\_refused\_at\_parse\_time}
\runs{tests/pql/test\_custom\_sql.py::test\_it\_runs\_on\_an\_engine\_it\_names\_and\_is\_refused\_on\_others}

\begin{definition}[Delegate]\label{def:delegate}
A \defn{delegate} is a registered Python class, named in PQL as
\pql{CHECK ds USING DELEGATE 'acme.settlement\_cycle@1' (\ldots) BELOW 0.1\%}, that returns
\emph{counts} --- rows scanned and rows violating --- and never a verdict. The control's threshold judges the
counts. A delegate's source hash is folded into the evidence.
\end{definition}

\begin{proposition}[A delegate cannot decide]\label{prop:delegate}
The verdict of a delegate control is $\mathrm{judge}(m, \theta)$ of \Cref{def:verdict} applied to the
delegate's counts, so a delegate that establishes nothing yields $\IND$, not $\PASS$. A delegate importing a
network, clock or model client is refused before it is imported; one that is not deterministic on admission
probes is refused; a file changed after admission is refused; input over the row ceiling is refused, not
truncated; and a sandboxed run agrees with an in-process run.
\end{proposition}
\runs{tests/delegates/test\_delegates.py::test\_a\_result\_that\_establishes\_nothing\_is\_indeterminate\_not\_a\_pass}
\runs{tests/delegates/test\_delegates.py::test\_an\_impure\_file\_is\_refused\_before\_it\_is\_imported}
\runs{tests/delegates/test\_delegates.py::test\_admission\_refuses\_what\_cannot\_replay\_or\_be\_judged}
\runs{tests/delegates/test\_streaming\_and\_testkit.py::test\_a\_file\_changed\_after\_admission\_is\_refused\_in\_the\_sandbox}
\runs{tests/delegates/test\_delegates.py::test\_too\_many\_rows\_are\_refused\_not\_truncated}
\runs{tests/delegates/test\_delegates.py::test\_sandboxed\_and\_in\_process\_runs\_agree\_and\_both\_see\_json\_values}

The sandbox is a subprocess with CPU and memory limits and one wall-clock deadline over both of its pipes.
Its environment is an allowlist, so the database's DSN and a provider's key are not in it. Where the host
permits unprivileged namespaces it runs in a network namespace of its own, with no interface up. In every
case an audit hook (PEP 578), installed before the delegate is imported and not removable afterwards, refuses
sockets, processes, \texttt{exec}, \texttt{fork}, and \texttt{ctypes}. The evidence records which of these
applied. Each is exercised by a delegate that passes admission and misbehaves only when a control asks it to.
Two findings came out of building the tests. The deadline had covered only the reply, so a delegate that
slept (using no CPU, so the CPU limit never fired) held the host for ever. And the admission scan let
\texttt{getattr(len.\_\_self\_\_, "\_\_imp" + "ort\_\_")} through; it now refuses \texttt{builtins}, the
interpreter-internals attributes and \texttt{getattr} with a computed name. Admission is sandboxed as well:
the server never imports delegate code. Each configured file is scanned, imported and probed by its own
sealed worker, which returns a description and the implementation hash, and a delegate that hangs its
admission is refused alone. One limit remains: an audit hook is not a boundary against native code already
loaded.
\runs{tests/delegates/test\_admission.py::test\_the\_server\_never\_imports\_a\_configured\_delegate}
\runs{tests/delegates/test\_admission.py::test\_a\_delegate\_that\_hangs\_its\_admission\_is\_refused\_alone}
\runs{tests/delegates/test\_sandbox.py::test\_what\_a\_delegate\_has\_no\_use\_for\_is\_refused}
\runs{tests/delegates/test\_sandbox.py::test\_a\_sleeping\_delegate\_is\_stopped\_at\_its\_deadline}
\runs{tests/delegates/test\_sandbox.py::test\_memory\_is\_limited}
\runs{tests/delegates/test\_sandbox.py::test\_the\_worker\_gets\_no\_secrets\_from\_the\_environment}

\begin{worked}
\wk{a check PQL cannot say}
Case study 5 plants 23 US trades booked on T+2 after the move to T+1, 9 EU trades settling on a TARGET2
closing day, and 4 unreadable dates. The column comparison \pql{settlement\_date >= trade\_date} passes all
of them, because every one settles after it trades. A bank-owned delegate that knows each market's cycle
and calendar returns 36 violations (23 late, 9 early, 4 unreadable); the threshold, not the delegate, turns
them into a failure. A second delegate, run on a remote agent inside a payments zone, applies Benford's
first-digit test and fails the planted ledger (MAD 0.0176, above Nigrini's 0.015 nonconformity band) while
keeping the example rows inside the zone.
\end{worked}

\subsection{Numbers}

Operator arithmetic in the function library and the reference interpreter is exact decimal: one tenth plus
two tenths is three tenths, and division by zero is unknown rather than an error. Threshold values are
parsed and stored as binary floating point, and survive rendering digit for digit; the comparison of a
metric with a threshold is therefore a float comparison.
\runs{tests/pql/test\_numbers\_survive.py::TestOperatorArithmeticIsExact::test\_a\_tenth\_plus\_two\_tenths\_is\_three\_tenths}
\runs{tests/pql/test\_numbers\_survive.py::TestRenderingKeepsEveryDigit::test\_a\_threshold\_survives\_a\_round\_trip}
\partly{thresholds are \texttt{float}, not decimal.}

% =============================================================================
\section{Evidence that verifies without its author}\label{sec:evidence}

\subsection{The record}

\begin{definition}[Evidence record and chain]\label{def:record}
An \defn{evidence record} $e_i$ is a mapping of fields --- tenant, control identity and version, plan
identity, snapshot, verdict, metrics, dimensions, a reference by hash to any sample rows, the format version,
the sequence number $i$ --- together with three hashes. Its \defn{content hash} is
$h_i = \hash(\mathrm{canon}(e_i \setminus \{\mathit{prev}, h, \mathit{rh}\}))$, computed over exactly the
fields its own format version defines. Its \defn{record hash} is $\mathit{rh}_i = \hash(\mathit{rh}_{i-1}
\,\Vert\, h_i)$, with $\mathit{rh}_{0} = \textsf{GENESIS} = \texttt{0}^{64}$. $\hash$ is SHA-256 and
$\mathrm{canon}$ is canonical JSON (keys sorted at every level, no whitespace, UTF-8). A chain is valid if
its sequence numbers are contiguous from 1, every $h_i$ recomputes, and every $\mathit{rh}_i$ links.
\end{definition}

\begin{theorem}[Tamper evidence]\label{thm:tamper}
Suppose a verifier holds an authentic chain head $\mathit{rh}_n$. If $\hash$ is collision resistant, then any
alteration of a field of any $e_i$ ($i \le n$), any deletion, insertion or reordering of records, or any
truncation of the chain, is detected by the validity check, except with negligible probability.
\end{theorem}
\begin{proof}[Proof sketch]
An altered field changes $\mathrm{canon}(e_i)$ and hence, except on a collision, $h_i$; if $h_i$ is left
unchanged the recomputation fails. If $h_i$ is recomputed, $\mathit{rh}_i$ changes, and by induction every
subsequent record hash changes, so $\mathit{rh}_n$ differs from the authentic head unless a collision was
found at some step. Deletion, insertion and reordering break contiguity of sequence numbers or the link
equation at the point of change, and a truncated chain ends at a head different from $\mathit{rh}_n$.
\end{proof}
\runs{tests/evidence/test\_ledger.py::TestTheChain::test\_a\_fresh\_chain\_starts\_from\_genesis}
\runs{tests/evidence/test\_ledger.py::TestTheChain::test\_each\_record\_links\_to\_the\_one\_before}
\runs{tests/evidence/test\_ledger.py::TestTheChain::test\_a\_missing\_record\_is\_detected\_as\_a\_gap}
\runs{tests/evidence/test\_ledger.py::TestTheChain::test\_reordering\_is\_detected}
\runs{tests/evidence/test\_independent\_verifier.py::TestItRejectsTampering::test\_a\_truncated\_file\_is\_caught}
\runs{tests/evidence/test\_independent\_verifier.py::TestItRejectsTampering::test\_a\_removed\_middle\_record\_is\_caught}
\runs{tests/evidence/test\_independent\_verifier.py::TestItRejectsTampering::test\_a\_reordered\_pair\_is\_caught}

The theorem is only as strong as the authenticity of the head. \textsc{Prama} signs the head with
HMAC-SHA256, which proves nothing to a party that does not hold the key --- the code says so in as many
words --- and seals exported bundles with an HMAC and, optionally, an Ed25519 signature over the manifest
that verifies with the public key alone. A bundle without a signature is reported untrustworthy even when
every hash matches.
\runs{tests/evidence/test\_ledger.py::TestSigning::test\_a\_different\_key\_does\_not\_verify}
\runs{tests/evidence/test\_ledger.py::TestSigning::test\_a\_changed\_head\_does\_not\_verify}
\runs{tests/security/test\_bundle.py::TestAsymmetricSignature::test\_a\_signature\_verifies\_with\_only\_the\_public\_half}
\runs{tests/security/test\_bundle.py::TestAnUnsignedBundleIsNotABundle::test\_no\_signature\_is\_not\_trustworthy\_even\_when\_the\_hashes\_match}
\paragraph{External anchoring.} Signing does not stop the signer. An operator holding the HMAC key can
rebuild a chain from genesis and re-sign its head, and the rebuilt chain verifies perfectly. What stops that
is a witness outside the operator's control \cite{haber1991,laurie2013}. After a run, \textsc{Prama} sends the
head's record hash (32 bytes, naming no record) to an RFC 3161 time-stamp authority, in a unit of work of its
own, so a witness that is down cannot roll the run's evidence back. It stores the token beside the chain. A
failed attempt is stored as such, not dropped. The offline verifier checks each receipt against the record at
its position, reading the token with the standard library alone. Given the authority's certificate, it
verifies the signature with \texttt{openssl ts -verify}, a tool the auditor already trusts. A rebuilt chain
then disagrees with a receipt its author cannot forge. What remains open is the window after the last anchor,
which is why anchoring happens after every run and not on a daily schedule.
\runs{tests/evidence/test\_anchor.py::test\_an\_auditor\_verifies\_the\_anchor\_and\_catches\_a\_rebuilt\_chain}
\runs{tests/evidence/test\_anchor.py::test\_the\_head\_is\_anchored\_once\_and\_a\_failure\_is\_recorded}
\runs{tests/evidence/test\_anchor.py::test\_residency\_is\_consulted\_before\_anything\_leaves}

\begin{definition}[Merkle root]\label{def:merkle}
For record hashes $x_1,\ldots,x_n$, the root is computed level by level, combining adjacent pairs with
$\hash(x_{2j-1} \Vert x_{2j})$; an unpaired node at the end of a level is \emph{promoted} unchanged to the
next level; the root of the empty set is \textsf{GENESIS} \cite{merkle1987}.
\end{definition}

\begin{remark}
Promotion rather than duplication matters: the common alternative, duplicating the last node, gives the
lists $[a,b,c]$ and $[a,b,c,c]$ the same root, which lets a duplicated record pass unnoticed.
\end{remark}
\runs{tests/evidence/test\_ledger.py::TestMerkleRoots::test\_an\_odd\_node\_is\_promoted\_not\_duplicated}
\runs{tests/evidence/test\_ledger.py::TestMerkleRoots::test\_an\_empty\_set\_has\_no\_root}

\subsection{A verifier that shares nothing}

\begin{proposition}[Independent verification]\label{prop:independent}
The script \texttt{scripts/verify\_evidence.py} imports only the Python standard library, runs where
\textsc{Prama} is not importable, and checks content hashes, links, contiguity, erasure seals, the manifest's
record count and file digest, the Merkle root and the chain head. On every input tested, including records
with non-ASCII text, an added field, and doctored manifests, it reaches the same conclusion as \textsc{Prama}'s
own verifier. A second, test-local re-implementation written from the prose description of the algorithm
also agrees.
\end{proposition}
\runs{tests/evidence/test\_independent\_verifier.py::TestItIsGenuinelyIndependent::test\_it\_imports\_only\_the\_standard\_library}
\runs{tests/evidence/test\_independent\_verifier.py::TestItIsGenuinelyIndependent::test\_it\_runs\_without\_prama\_importable}
\runs{tests/evidence/test\_independent\_verifier.py::TestTheTwoVerifiersAgree::test\_a\_non\_ascii\_record\_verifies\_in\_both}
\runs{tests/evidence/test\_independent\_verifier.py::TestTheTwoVerifiersAgree::test\_an\_added\_field\_is\_a\_breach\_in\_both}
\runs{tests/evidence/test\_independent\_verifier.py::TestItRejectsTampering::test\_a\_doctored\_manifest\_is\_caught}
\runs{tests/evidence/test\_ledger.py::TestVerifiableWithoutPrama::test\_an\_independent\_reader\_reaches\_the\_same\_conclusion}
\runs{tests/evidence/test\_ledger.py::TestVerifiableWithoutPrama::test\_and\_the\_same\_conclusion\_when\_a\_record\_is\_altered}
\runs{tests/evidence/test\_end\_to\_end.py::TestVerifyingWithoutPrama::test\_an\_exported\_chain\_checks\_with\_stdlib\_alone}

\begin{plain}
\pt An auditor does not have to trust \textsc{Prama} to check \textsc{Prama}'s records. A short script that
uses nothing but Python's standard library reads the exported file and says whether anything was changed,
removed or reordered. The two verifiers once disagreed (QA finding C4, since fixed); that disagreement was a
defect in the evidence layer, it was found by running the two against each other, and a test --- including a
record with non-ASCII text --- now holds them together.
\end{plain}

\paragraph{Erasure without breaking the chain.} A record that must be erased (a legal hold ends, or a subject
exercises a right to erasure) is replaced by a tombstone that keeps the original content hash, so the chain
still links, and carries a seal of its own over who erased it and why. Rewriting who erased a record is
caught.
\runs{tests/evidence/test\_independent\_verifier.py::TestAnErasureIsItselfTamperEvident::test\_rewriting\_who\_erased\_it\_is\_caught}
\runs{tests/evidence/test\_independent\_verifier.py::TestAnErasureIsItselfTamperEvident::test\_the\_chain\_still\_verifies\_across\_the\_erasure}

\paragraph{Format versions.} A field added to the record format is hashed only in records whose version
defines it, so a chain mixing old and new records verifies.
\runs{tests/evidence/test\_format\_versions.py::TestOldRecordsSurviveANewBuild::test\_a\_mixed\_chain\_verifies}

\paragraph{Size.} Sample rows are referenced by hash rather than carried. The tests assert that the median
record is under 2 KB; the roadmap reports 848 bytes on the worked example, a figure we did not reproduce.
\runs{tests/evidence/test\_ledger.py::TestSize::test\_a\_record\_is\_well\_under\_two\_kilobytes}
\runs{tests/evidence/test\_ledger.py::TestSize::test\_samples\_are\_referenced\_not\_carried}

\subsection{Replay, and naming why an answer moved}

\begin{definition}[Replay outcome]\label{def:replay}
Re-executing the plan of a record yields one of: \textsf{identical}; \textsf{stable} (the answer held while
the inputs moved); one of six named causes of divergence --- \textsf{data\_changed},
\textsf{control\_changed}, \textsf{engine\_changed}, \textsf{parameters\_changed},
\textsf{snapshot\_not\_exact}, \textsf{coverage\_changed}; or \textsf{unexplained}. Causes are tried in a fixed
precedence (a changed control explains more than changed data), and \textsf{unexplained} is never folded into
an ordinary cause.
\end{definition}

\begin{proposition}[Replay classifies]\label{prop:replay}
Each named cause is assigned when, and only when, its condition holds in the tested scenarios; a changed
control takes precedence over changed data; an engine disagreement is escalated; and a run whose answer held
is reported as \textsf{stable}, not diverged, so that a nightly replay against fresh data does not bury the
divergences that matter.
\end{proposition}
\runs{tests/evidence/test\_replay.py::TestNamingTheCause::test\_moved\_data\_is\_the\_ordinary\_case}
\runs{tests/evidence/test\_replay.py::TestNamingTheCause::test\_an\_edited\_control\_makes\_the\_answers\_incomparable}
\runs{tests/evidence/test\_replay.py::TestNamingTheCause::test\_the\_control\_wins\_when\_both\_the\_control\_and\_the\_data\_moved}
\runs{tests/evidence/test\_replay.py::TestNamingTheCause::test\_an\_engine\_disagreement\_is\_escalated}
\runs{tests/evidence/test\_replay.py::TestNamingTheCause::test\_an\_inexact\_snapshot\_explains\_itself}
\runs{tests/evidence/test\_replay.py::TestNamingTheCause::test\_an\_unexplained\_divergence\_is\_reported\_as\_alarming}
\runs{tests/execute/test\_watermarks.py::TestTheWidthTravelsWithTheVerdict::test\_a\_replay\_at\_a\_different\_width\_is\_not\_blamed\_on\_the\_data}
\runs{tests/evidence/test\_end\_to\_end.py::TestReplay::test\_the\_controls\_whose\_answer\_held\_are\_not\_reported\_as\_diverged}

Every case study ends by verifying its chain and printing its Merkle root; all eight verified in our runs,
over 10 to 81 records each.

% =============================================================================
\section{Reconciliation}\label{sec:recon}

Reconciliation --- showing that two systems that should agree do --- is the most expensive data quality
failure in banking, and is usually sold as a separate product. In \textsc{Prama} it is a PQL statement:
\begin{verbatim}
RECONCILE subledger AGAINST general_ledger
  ON (account, cost_centre, posting_date)
  COMPARING amount = balance_eur WITHIN 0.01 EUR
  NORMALISING currency TO 'EUR' USING RATES fx_rates
  OFFSET BY 1 DAY
\end{verbatim}

\begin{definition}[Reconciliation]\label{def:recon}
A reconciliation $\mathcal{R} = (L, R, M, (c_L, c_R), \epsilon, \nu, k)$ has a left and a right dataset, a
match key $M$, compared attributes, a tolerance $\epsilon = (a, r)$ with absolute bound $a \ge 0$ and
relative bound $r \ge 0$ (either may be absent, treated as $0$), an optional normalisation $\nu$ converting
the left amount into the right's currency through a declared rate table, and an optional offset of
$k$ days on one date component of the key.
\end{definition}

\begin{definition}[Tolerance]\label{def:tolerance}
A difference $\Delta$ against a magnitude $m$ is \defn{permitted} when
$|\Delta| \le a \;\lor\; |\Delta| \le r\,|m|$ --- ``whichever is larger''. A pair is a \emph{break} when
the difference is not permitted, that is, when it breaches \emph{both} bounds.
\end{definition}

\begin{proposition}[Tolerance laws]\label{prop:tolerance}
(i) Permission is monotone: enlarging $a$ or $r$ never creates a break. (ii) For large positions the relative
bound governs; for small ones the absolute bound. (iii) Against a zero magnitude, a purely relative bound
permits only an exact match. (iv) A difference breaching both bounds is a break however large either bound is
individually.
\end{proposition}
\begin{proof}
(i) Each disjunct is monotone in its bound. (ii) $r|m| > a$ exactly when $|m| > a/r$. (iii) $r\cdot 0 = 0$.
(iv) is \Cref{def:tolerance} restated.
\end{proof}
\runs{tests/recon/test\_reconciliation.py::TestWhicheverIsLarger::test\_the\_relative\_bound\_governs\_a\_large\_position}
\runs{tests/recon/test\_reconciliation.py::TestWhicheverIsLarger::test\_a\_difference\_breaching\_both\_is\_still\_a\_break}
\runs{tests/recon/test\_reconciliation.py::TestWhicheverIsLarger::test\_a\_relative\_bound\_against\_nothing\_permits\_only\_an\_exact\_match}
\runs{tests/recon/test\_reconcile\_pql.py::test\_the\_engine\_counts\_genuine\_breaks\_and\_the\_threshold\_judges}
\runs{tests/recon/test\_reconcile\_pql.py::test\_normalising\_converts\_before\_comparing}

\begin{definition}[Matching with an offset]\label{def:offsetmatch}
Rows are first matched exactly on $M$. Only left rows still unmatched are then probed at offsets
$\pm 1, \ldots, \pm k$ days on the named date component of $M$ (by default the last), against right rows
also still unmatched.
\end{definition}

\begin{lemma}[An offset never steals an exact match]\label{lem:exactfirst}
No row that has an exact match is paired across an offset.
\end{lemma}
\begin{proof}Exact matching runs to completion before any offset probe, and probes consider only rows left
unmatched by it.\end{proof}
\runs{tests/recon/test\_reconciliation.py::test\_an\_exact\_match\_is\_never\_quietly\_paired\_with\_the\_wrong\_day}
\runs{tests/recon/test\_reconciliation.py::test\_a\_date\_one\_day\_apart\_is\_matched\_rather\_than\_double\_counted}

\begin{proposition}[The price of ignoring timing]\label{prop:timing}
Let $T$ be a set of keys whose right-hand entry is posted $\ell \le k$ days late, such that for each
$t \in T$ neither the original date nor the late date has another unmatched counterpart and the amounts agree
within tolerance. Let $B_k$ be the number of breaks with offset $k$ and $B_0$ the number without. Then
\[ B_0 = B_k + 2|T|. \]
\end{proposition}
\begin{proof}
With the offset, by \Cref{lem:exactfirst} and the hypothesis, each $t \in T$ is paired across the offset
and contributes no break. Without it, $t$ has a left row at date $d$ with no right row at $d$ (a
\emph{missing}) and a right row at $d+\ell$ with no left row there (an \emph{extra}): two breaks. Breaks not
involving $T$ are unaffected, since exact matching is the same in both runs.
\end{proof}
\runsprose{case study 6, run by us: 5 breaks with \pql{OFFSET BY 1 DAY}, 9 without, with $|T| = 2$.
No test asserts either number; \texttt{run.py} prints them and \texttt{README.md} states them.}

\begin{worked}
\wk{a month-end close}
Case study 6 builds an ERP with 840 subledger rows, 280 ledger rows and three currencies, and plants four
defects: a manual journal of 250.00 EUR posted to the ledger only, a day's entries that never reached the
ledger, a ledger balance on an account the subledger never booked, and two entries from the last evening
that reached the ledger the next morning. The reconciliation above found five breaks needing a person ---
three \emph{genuine} (the journal, on three dates), one \emph{missing}, one \emph{extra} --- and no break for
the timing difference. Run without \pql{OFFSET BY 1 DAY}, the same books show nine: each of the two late
entries becomes one missing on the last day and one extra the next, exactly as
\Cref{prop:timing} predicts.
\end{worked}

\begin{definition}[Break classes]\label{def:breaks}
Each break is classified as \textsf{timing}, \textsf{fx}, \textsf{rounding}, \textsf{missing},
\textsf{extra}, \textsf{duplicate}, \textsf{sign} or \textsf{genuine}. Only \textsf{timing} is expected to
clear by itself; \textsf{sign}, \textsf{duplicate} and \textsf{fx} are configuration faults and are routed away
from the data steward; the violation count of a reconciliation is the number of breaks that need a person.
\end{definition}
\runs{tests/recon/test\_reconciliation.py::test\_a\_timing\_difference\_is\_expected\_to\_clear}
\runs{tests/recon/test\_reconciliation.py::test\_a\_sign\_convention\_error\_is\_named\_rather\_than\_reported\_as\_a\_big\_break}
\runs{tests/recon/test\_reconciliation.py::test\_a\_duplicated\_posting\_is\_named}
\runs{tests/recon/test\_reconciliation.py::test\_configuration\_faults\_are\_routed\_away\_from\_the\_data\_steward}

The statement round-trips through the parser and renders as a sentence. A filter is written after the
dataset it applies to, \pql{RECONCILE a WHERE p AGAINST b WHERE q}, and a single \pql{WHERE} on the whole
control is refused: a filter on one side only reports every row it dropped from that side as missing on the
other. An unfiltered reconciliation's plan is unchanged by the addition, so evidence already written still
names it.
\runs{tests/recon/test\_reconcile\_pql.py::test\_it\_round\_trips\_and\_says\_what\_it\_checks}
\runs{tests/recon/test\_reconcile\_pql.py::test\_a\_filter\_on\_the\_whole\_control\_is\_refused}
\runs{tests/recon/test\_reconcile\_pql.py::test\_each\_side\_filters\_its\_own\_rows}
\runs{tests/lineage/test\_joins.py::test\_a\_filtered\_copy\_reconciles\_against\_the\_same\_rows\_of\_its\_source}
\absent{user-written \pql{CLASSIFY} rules on \pql{RECONCILE} (the engine's built-in classification is used);
a signed reconciliation certificate from the break workbench.}

% =============================================================================
\section{Lineage read from code}\label{sec:lineage}

\subsection{Parsed, never executed}

\begin{definition}[Lineage graph]\label{def:lineage}
A \defn{column lineage graph} is a directed graph whose nodes are qualified columns and whose edges carry a
transform $\tau$, one of \textsf{identity}, \textsf{rename}, \textsf{derived}, \textsf{aggregated},
\textsf{filter} and \textsf{join\_key}; a status, \textsf{parsed} or \textsf{inferred}; a confidence;
and the source span it was read from.
\end{definition}

\textsc{Prama} reads SQL with sqlglot, falling back to a pattern reader that records the fallback as a gap;
PySpark, pandas and Airflow code are read with Python's \texttt{ast} module (importing a DAG file runs it, so
it is never imported); Power BI models are read as JSON. An edge produced by a parser is \textsf{parsed} with
confidence 1.0; an edge produced by the fallback is \textsf{inferred} at 0.8; an edge proposed by a language
model is \textsf{inferred} and capped at 0.85, and is kept only if it is checked against the code.
\runs{tests/lineage/test\_pyspark.py::test\_the\_code\_is\_parsed\_never\_run}
\runs{tests/codeintake/test\_intake.py::test\_nothing\_received\_is\_executable}
\runs{tests/lineage/test\_sql.py::test\_the\_parser\_path\_is\_the\_one\_running}
\runs{tests/lineage/test\_sql.py::test\_what\_the\_parser\_cannot\_read\_falls\_back\_and\_says\_so}
\runs{tests/lineage/test\_sql.py::test\_a\_where\_clause\_column\_feeds\_the\_target\_as\_a\_filter}
\runs{tests/codeintake/test\_model\_lineage.py::test\_the\_model\_pass\_keeps\_only\_checked\_edges\_as\_inferred}

On a gold fixture ETL repository the extractor is required to reach precision at least 0.98 and recall at
least 0.95 on value edges, with filter and join-key edges excluded from the measure.
\runs{tests/codeintake/test\_gold\_fixture.py::test\_bankco\_etl\_meets\_the\_gate}

\subsection{Blast radius}

\begin{definition}[Impact]\label{def:impact}
Each transform has an attenuation $\alpha(\tau) \in (0,1]$: identity and rename $1.0$, filter $0.9$,
join key $0.8$, derived $0.7$, aggregated $0.35$. The impact of a path is the product of its attenuations,
and the \defn{impact} of a defect at $s$ on a node $v$ is the maximum over paths from $s$ to $v$ of length at
most 12, reported when it is at least $0.01$.
\end{definition}

\begin{proposition}[Strongest path]\label{prop:strongest}
Because every $\alpha \le 1$, extending a path never increases its impact. A breadth-first search that
re-expands a node whenever a strictly stronger path reaches it (a label-correcting search) therefore
terminates, even on a cyclic graph, with every node at its maximum-product impact; the first path found is not
in general the strongest.
\end{proposition}
\runs{tests/lineage/test\_graph.py::test\_the\_strongest\_path\_wins\_rather\_than\_the\_first\_one\_found}
\runs{tests/lineage/test\_graph.py::test\_an\_aggregate\_attenuates\_more\_than\_a\_copy}
\runs{tests/lineage/test\_graph.py::test\_a\_filter\_is\_barely\_attenuated}
\runs{tests/lineage/test\_graph.py::test\_a\_negligible\_impact\_is\_counted\_rather\_than\_listed}

The attenuations are declared constants, not estimates. They order transforms sensibly and they carry no
statistical meaning; a reader should read ``35\% of the defect'' as a ranking, not a measurement. The code
also contains a comment saying a filter ``is not attenuated at all'' beside the value 0.9; the value is what
runs.
\stated{the attenuation constants are not calibrated against observed defect propagation.}

\subsection{Controls justified by lineage}

A lineage edge is a claim about how one column is computed from another. Some controls follow from that
claim, and some do not.

\begin{definition}[Lineage-derived proposals]\label{def:linprop}
From a lineage graph and the active controls, \textsc{Prama} proposes: \emph{propagation}, carrying a
row-level control on $s$ to $t$ when $t$ is an identity or rename copy of $s$; \emph{referential}, a
\pql{REFERENCES} check from a copied key column to its source; and \emph{reconciliation}, a
\pql{RECONCILE t AGAINST s \ldots\ WITHIN 0} when keys and an amount are copied at the same grain. A
proposal is \emph{held}, with its reason, when any edge it rests on is \textsf{inferred}, or --- for a
reconciliation --- when a \textsf{filter} edge lies between source and target.
\end{definition}

\begin{proposition}[What a copy justifies]\label{prop:copy}
Let $t$ be produced from $s$ by selecting a subset of rows and copying a column by identity or rename.
(i) For any row predicate $\varphi$, if every row of $s$ satisfies $\varphi$ then every row of $t$ does; so
propagating a row-level control from $s$ to $t$ is sound, filtered or not. (ii) A reconciliation of $t$
against $s$ on their keys reports every row the filter removed as missing, so it is sound only when no filter
lies between them. (iii) Neither (i) nor (ii) holds for an aggregated column.
\end{proposition}
\begin{proof}
(i) Every row of $t$ is a copy of a row of $s$. (ii) Rows of $s$ outside the filter have no counterpart in
$t$. (iii) An aggregate is not a copy of any row, so a row-level predicate on the input says nothing about
the output, and the output's grain differs.
\end{proof}
\runs{tests/derive/test\_lineage\_controls.py::test\_a\_control\_is\_carried\_to\_a\_copied\_column}
\runs{tests/derive/test\_lineage\_controls.py::test\_an\_aggregated\_column\_does\_not\_inherit\_a\_row\_level\_control}
\runs{tests/derive/test\_lineage\_controls.py::test\_a\_copied\_key\_must\_exist\_in\_its\_source}
\runs{tests/derive/test\_lineage\_controls.py::test\_a\_proposal\_on\_an\_inferred\_edge\_is\_held}
\runs{tests/derive/test\_lineage\_controls.py::test\_every\_proposal\_is\_valid\_pql}
\runs{tests/recon/test\_reconcile\_pql.py::test\_lineage\_proposes\_a\_reconciliation\_for\_a\_copy\_at\_the\_same\_grain}

\begin{plain}
\pt If staging is a straight copy of the raw feed, a rule that must hold on every raw row must hold on
every staged row too, so \textsc{Prama} proposes the same rule downstream. But if staging keeps only booked
trades, reconciling it against the raw feed would report every cancelled trade as a break, so that
proposal is held back with the reason. A guess about lineage never becomes a control without somebody
confirming the guess first.
\end{plain}

Case study 8 found both halves of \Cref{prop:copy}(ii) the hard way: the first run proposed a
reconciliation across the filter and reported 118 breaks, one per cancelled trade, and the proposed
reconciliation had no tolerance and could not run. Both are fixed. The proposal says \pql{WITHIN 0}, and
the filter edge now carries its condition when it is over one table, so the proposal applies the same filter
to the source side (\pql{AGAINST raw.trades WHERE status = 'BOOKED'}); a condition that cannot be carried over
holds the proposal, with the reason.
\runs{tests/recon/test\_reconcile\_pql.py::test\_lineage\_proposes\_a\_reconciliation\_for\_a\_copy\_at\_the\_same\_grain}

\subsection{A negative result: value lineage cannot see a join}\label{sec:join}

\begin{proposition}[Join predicates are invisible to value lineage]\label{prop:join}
Let $m$ be a column of a view $Q = \pi_{\ldots m \ldots}(S \bowtie_{S.c = F.c} F)$ and suppose $c$ is not
projected into $Q$. In a lineage graph whose edges record only value flow (identity, rename, derived,
aggregated) and filters on \pql{WHERE} columns, there is no path from $S.c$ to $Q.m$; yet a defect in $S.c$
that makes the join predicate false removes rows from $Q$, and so changes every aggregate of $m$.
\end{proposition}
\begin{proof}
No value of $S.c$ is copied, transformed or aggregated into any column of $Q$, and $c$ is not in a
\pql{WHERE} clause, so the extractor emits no edge from $S.c$ into $Q$. The dependence of $Q$ on $S.c$ is a
dependence of $Q$'s \emph{row population}, not of any column's values.
\end{proof}
\runsprose{case study 8, run by us. No test asserts it; the nearest test is
\texttt{tests/lineage/test\_pyspark.py::test\_a\_join\_is\_a\_gap\_not\_a\_guess}, which requires the PySpark
reader to record a join as a gap rather than invent an edge.}

\begin{worked}
\wk{four lower-case currencies, 605 million of notional}
Case study 8 builds a warehouse by running an ETL repository: raw feed, staging (booked trades only), and a
mart that joins staging to an FX table on currency and aggregates exposure, with a Power BI model on top.
Intake read three files and, when the study was first built, emitted twelve parsed column edges. Two defects are planted in the raw feed: five
notionals carrying the sign of the side, and four currencies in lower case. The owner writes two controls on
the raw columns; lineage proposes the rest. The notional's blast radius runs to the dashboard --- 100\% at
staging, 35\% at the mart's exposure (aggregated), 35\% at the dashboard's exposure (renamed), 12\% at its
total ($0.35^2$) --- and the propagated control finds the five at staging too. The currency's blast radius
\emph{stops at staging}: the mart uses currency only in the join. Yet its effect there is real, and the run
measures it: three staged trades with no FX rate drop out of the mart, 605,000,000 of notional with them, with
no error, no null and no count anywhere that changes. Only the control on the raw column saw it.
\end{worked}

\paragraph{The repair: a population edge.} The proposition is about a graph of value flow, so the repair
is not a better value edge but a different kind of node. The SQL reader now records each equality in an
inner or left join's \pql{ON} clause as two \textsf{join\_key} edges, one from each side's key, into the
view's \emph{rows}, written $Q.\ast$ --- the same node a \pql{WHERE} column already fed --- and reads a
common table expression through to its real table. The blast radius then treats $Q.\ast$ as a population
node: a change in which rows exist reaches every column computed over them, and the path says so
(\emph{computed over these rows}). The edge does not pretend to be a value effect: nothing flows from $S.c$
into $Q.m$ except through $Q.\ast$, and a reader sees the rows named as the thing affected. What remains a
heuristic is its weight, $0.8$, which is declared rather than estimated like every other attenuation
(\Cref{sec:lineage}).
\runs{tests/lineage/test\_joins.py::test\_a\_join\_key\_is\_recorded\_through\_a\_cte\_to\_the\_real\_table}
\runs{tests/lineage/test\_joins.py::test\_a\_cte\_is\_not\_mistaken\_for\_a\_table}
\runs{tests/lineage/test\_joins.py::test\_the\_blast\_radius\_follows\_a\_join\_to\_the\_dashboard}

The join also justifies a control, by the same argument as \Cref{prop:copy}: in
$S \bowtie_{S.c = F.c} F$ every $S$ row whose $c$ has no match in $F$ is dropped by an inner join, or reaches
$Q$ with nothing joined to it by a left join, and in neither case does anything fail. So lineage proposes
\pql{CHECK S.c REFERENCES F.c DIMENSION integrity} on the driving side, held while the edge is inferred.
\runs{tests/recon/test\_reconcile\_pql.py::test\_a\_join\_proposes\_that\_every\_driving\_row\_finds\_its\_match}

Rerun with both changes, case study 8 reads fourteen edges. The proposed join check fails on 3 of 1,882
staged trades, which is where the defect happens and not only where it entered, and the currency's blast
radius now runs from staging to the mart's rows (80\%), the mart's exposure (80\%), the dashboard's exposure
(80\%) and its total (28\%). The filtered reconciliation of staging against the raw feed passes: the copy is
faithful, and the defects are in what it faithfully copied. We still take the result to mean that controls
belong at the source; lineage carries them downstream, and now carries them across a join.
The PySpark reader follows a join too: \texttt{on="k"}, a list of names, or an equality between the two
frames. It records the key pairs with the same pairing text, and takes a column after the join from the
side that provably has it. A name either side could hold is a gap, as in SQL.
\runs{tests/lineage/test\_pyspark.py::test\_a\_join\_is\_followed\_to\_its\_keys\_and\_its\_columns}
\runs{tests/lineage/test\_pyspark.py::test\_a\_pyspark\_join\_proposes\_the\_same\_check\_as\_a\_sql\_one}
\absent{join-key edges from the pandas reader, where a \texttt{merge} is still a gap; population edges from
\pql{EXISTS} and \pql{IN} subqueries.}

\paragraph{Before it merges.} The same comparison run on a pull request answers the reviewer's questions.
\texttt{prama code review} reads the base and the head from git, never checking either out, through the
same sandboxed reader, and reports four things:
\begin{itemize}[nosep]
  \item the edges the change adds, removes or retypes (a copy becoming a derivation is a change);
  \item the blast radius of each changed column;
  \item the proposals new at the head;
  \item the proposals present at the base and gone at the head.
\end{itemize}
A live control whose identity is among the last rests on lineage the change removed. The review names it
and exits 3, so continuous integration can refuse the merge.
\runs{tests/codeintake/test\_review.py::test\_a\_live\_control\_that\_loses\_its\_basis\_fails\_the\_review}
\runs{tests/codeintake/test\_review.py::test\_a\_copy\_that\_becomes\_a\_derivation\_is\_a\_change\_with\_reach}

% =============================================================================
\section{Models author; the engine decides}\label{sec:ai}

\subsection{The boundary}

Language models are good at drafting a rule from a data dictionary, at ranking datasets for a question put
in plain words, and at explaining a failure. They are not accountable for a verdict, and their output is
not reproducible in the sense evidence requires. \textsc{Prama}'s rule is that a model may \emph{author},
\emph{rank}, \emph{explain}, \emph{calibrate} and \emph{summarise}, and that no code path lets a model's output
determine whether data passed.

\begin{definition}[Model-free verdict path]\label{def:modelfree}
The \defn{verdict path} is the sequence parse $\to$ type-check $\to$ lower $\to$ compile $\to$ execute
$\to$ judge $\to$ record. It is \defn{model-free} if no function on it takes as input a value returned by a
model client.
\end{definition}

\begin{proposition}[What the architectural guard establishes]\label{prop:guard}
The architecture test fails the build if any source module both refers to a model (by names such as
\texttt{llm}, \texttt{openai}, \texttt{completion}, \texttt{build\_prompt}) and produces or assigns a verdict
(\texttt{verdict}, \texttt{pass\_fail}, \texttt{is\_violation}, \texttt{assert\_outcome}), with docstrings
stripped before scanning. A counterfactual module that calls a model and returns a verdict makes the guard
fire. The guard is a syntactic necessary condition for \Cref{def:modelfree}, not a proof of it: a model's
output laundered through a neutral name in a separate module would not be caught.
\end{proposition}
\runs{tests/architecture/test\_layering.py::TestModelVerdicts::test\_no\_module\_both\_calls\_a\_model\_and\_produces\_a\_verdict}
\runs{tests/architecture/test\_layering.py::TestModelVerdicts::test\_the\_guard\_still\_fires\_on\_a\_module\_that\_would\_break\_the\_rule}
\runs{tests/architecture/test\_layering.py::TestModelVerdicts::test\_the\_guard\_ignores\_prose\_about\_the\_rule}
\runs{tests/architecture/test\_layering.py::TestModelVerdicts::test\_the\_exemption\_is\_a\_ban\_list\_not\_a\_model\_call}
\stated{an information-flow proof of \Cref{def:modelfree}; the guard approximates it by co-occurrence.}

The structural argument is stronger than the guard. A verdict is $\mathrm{judge}(m, \theta)$
(\Cref{def:verdict}), where $m$ comes from an engine executing a compiled plan and $\theta$ from the plan
itself; a plan comes from parsing PQL text that a person activated (\Cref{sec:proposals}). A model can
influence a verdict only by authoring PQL that a person then activates, at which point the verdict is a
function of that person's decision and the data, and the evidence names both. Validators and delegates that
import a model client are refused at admission (\Cref{prop:delegate}).

\subsection{Induced controls must be able to fail}

A model asked for a rule can return one that is well formed, type-correct, compiles, runs --- and cannot
fail: \pql{CHECK trades.side MATCHES /.*/}, or a bound of \pql{qty >= -1e308}.

\begin{definition}[Validation gates]\label{def:gates}
An induced control is \emph{validated} only after five gates: it parses; it type-checks against the
catalogue; it compiles; it runs in a sandbox over at most 2,000 sample rows; and a \emph{counterfactual} gate
finds data on which it fails. The type \texttt{Validated} cannot be constructed unless all five passed.
\end{definition}

\begin{proposition}[No vacuous control reaches a reviewer]\label{prop:vacuous}
A control that accepts every value it could be shown is rejected at the counterfactual gate, however well
formed, and never enters the proposal queue.
\end{proposition}
\runs{tests/induce/test\_validate.py::test\_a\_validated\_control\_cannot\_be\_built\_without\_passing\_every\_gate}
\runs{tests/induce/test\_validate.py::test\_a\_control\_that\_cannot\_fail\_is\_rejected\_however\_well\_formed}
\runs{tests/induce/test\_validate.py::test\_a\_bound\_so\_wide\_that\_nothing\_breaches\_it\_is\_rejected}
\runs{tests/induce/test\_inducer.py::test\_a\_vacuous\_answer\_never\_reaches\_a\_reviewer}

\begin{plain}
\pt A check that cannot fail is worth nothing, and a model is quite capable of writing one that looks
reasonable. So before anybody is shown a model's suggestion, \textsc{Prama} tries to make it fail. If it cannot,
the suggestion is thrown away.
\end{plain}

\subsection{The gateway}

All model calls go through one gateway. A caller asks for a \emph{purpose} (author, explain, discover,
embed), never for a model; a profile maps each purpose to an ordered route of providers. A purpose with no
profile is routed to a mock provider that answers with nothing --- so a fresh installation uses no model at
all, and still records the call.

\begin{definition}[Call ledger]\label{def:callledger}
Every call is appended to a ledger whose entries are sealed as in \Cref{def:record}: each seal is a hash over
the entry and the previous seal. Payloads are redacted by default (secrets, card numbers, IBANs) and may
expire; an expired payload is blanked and the chain still verifies.
\end{definition}
\runs{tests/llm/test\_gateway.py::test\_the\_seal\_chains}
\runs{tests/llm/test\_governance.py::test\_an\_expired\_payload\_is\_blanked\_and\_the\_chain\_still\_verifies}
\runs{tests/llm/test\_governance.py::test\_payloads\_follow\_policy\_and\_are\_redacted\_by\_default}
\runs{tests/llm/test\_gateway.py::test\_an\_unconfigured\_purpose\_falls\_to\_the\_mock\_which\_answers\_nothing}
\partly{\texttt{prama llm verify} is not itself exercised by a test; the method it calls is.}

\begin{proposition}[Evaluation gates activation]\label{prop:evalgate}
A prompt template version is approved only after a passing evaluation run, and not by its author; a new
profile version cannot be activated without its own passing run; and with no model configured no evaluation
can pass. Graders are deterministic.
\end{proposition}
\runs{tests/llm/test\_governance.py::test\_a\_template\_version\_is\_approved\_only\_after\_a\_passing\_run\_and\_by\_another}
\runs{tests/llm/test\_governance.py::test\_with\_the\_gate\_a\_new\_profile\_version\_waits\_for\_its\_own\_passing\_run}
\runs{tests/llm/test\_governance.py::test\_with\_no\_model\_configured\_nothing\_can\_pass}

Personal data may not be sent to a hosted model; when no provider on the route may receive a payload, the
refusal is recorded and raised; a spent budget refuses with HTTP 429.
\runs{tests/llm/test\_providers.py::test\_pii\_may\_not\_be\_sent\_to\_a\_hosted\_model}
\runs{tests/llm/test\_gateway.py::test\_when\_nothing\_may\_receive\_it\_the\_refusal\_is\_recorded\_and\_raised}
\runs{tests/llm/test\_providers.py::TestTheThreeDefects::test\_a\_secret\_and\_a\_card\_number\_are\_withheld\_from\_the\_prompt}
\runs{tests/llm/test\_providers.py::TestPatternRedaction::test\_a\_valid\_iban\_is\_withheld\_and\_a\_lookalike\_is\_not}
\runs{tests/api/test\_llm\_api.py::test\_a\_spent\_budget\_refuses\_with\_429}

\absent{any measurement of the \emph{quality} of model-authored rules --- acceptance rate, precision, or the
comparison of mining, induction and their fusion. The mechanism for merging corroborated origins runs
(\Cref{prop:corrob}); the claim that fusion beats either alone is about reviewers' decisions and was not
measured.}

% =============================================================================
\section{Trust and false-alarm budgets}\label{sec:trust}

\subsection{Trust as a choice of semiring}

A score for a derived column should depend on the scores of what it is derived from. How it should depend is
a modelling choice, and the choice is algebraic.

\begin{definition}[Trust propagation]\label{def:trust}
Given a lineage DAG, base scores $s(v) \in [0,1]$, and a pair of operations $(\otimes, \oplus)$ for combining
along a path and across paths, the propagated trust of $v$ combines its own score with
$\bigoplus_{\text{paths } p \to v} \bigotimes_{u \in p} s(u)$. \textsc{Prama} offers four choices:
\begin{center}\small
\begin{tabular}{@{}lll@{}}
\toprule
\textbf{Name} & \textbf{along a path} ($\otimes$) & \textbf{across paths} ($\oplus$) \\
\midrule
all inputs matter (default) & product & min \\
redundant sources & product & max \\
weakest link & min & min \\
product--mean & product & mean \\
\bottomrule
\end{tabular}
\end{center}
Attenuation (\Cref{def:impact}) applies to the \emph{deficit} $1 - s$, not to the score: an aggregate
passes on only part of a source's defect.
\end{definition}

\begin{remark}
The code calls all four ``semirings''. Strictly, only one is. \emph{Redundant sources},
$([0,1],\max,\times,0,1)$, is the Viterbi semiring of the provenance literature \cite{green2007}. The
default, $(\min,\times)$, has both operations associative and commutative and $\times$ distributes over
$\min$ on $[0,1]$, but the identity of $\min$ is $1$, which does not annihilate under $\times$; the same
holds for \emph{weakest link}, $(\min,\min)$. \emph{Product--mean} is not associative across paths at all,
since a mean of means over unequal groups is not the mean. Nothing in \textsc{Prama} relies on annihilation
--- a node with no inbound path is simply scored on its own evidence --- so the gap is one of naming, and we
record it rather than repair it here. The choice is exposed rather than hidden because the choices genuinely
disagree on the same graph.
\end{remark}
\runs{tests/score/test\_trust.py::test\_the\_default\_treats\_inputs\_as\_complementary}
\runs{tests/score/test\_trust.py::test\_the\_semirings\_genuinely\_disagree}
\runs{tests/score/test\_trust.py::test\_a\_copy\_inherits\_the\_defect\_intact}
\runs{tests/score/test\_trust.py::test\_an\_aggregate\_makes\_a\_column\_more\_trustworthy\_than\_its\_source}
\runs{tests/score/test\_trust.py::test\_a\_control\_that\_caught\_the\_problem\_stops\_it\_propagating}
\stated{the algebraic laws (associativity, distributivity) of each choice; the tests check behaviour on
examples, not the laws.}

\subsection{Conformal $p$-values}

A statistical monitor raises an alert when an observation is unusual. ``Unusual'' needs a false-alarm rate,
and a rate needs a guarantee.

\begin{definition}[Conformal $p$-value]\label{def:conformal}
Given calibration scores $s_1, \ldots, s_n$ and a new score $s$, the conformal $p$-value is
\[ p(s) = \frac{1 + |\{ i : s_i \ge s \}|}{n + 1}. \]
An alert is raised at level $\alpha$ when $p(s) \le \alpha$.
\end{definition}

\begin{theorem}[Validity {\cite{vovk2005}}]\label{thm:conformal}
If $s_1, \ldots, s_n, s$ are exchangeable, then $\Pr[p(s) \le \alpha] \le \alpha$ for every $\alpha \in [0,1]$.
\end{theorem}
\begin{proof}[Proof sketch]
Under exchangeability the rank of $s$ among the $n+1$ scores is uniform over $\{1, \ldots, n+1\}$ (breaking
ties conservatively), and $p(s) \le \alpha$ requires $s$ to rank among the top $\lfloor \alpha(n+1) \rfloor$.
\end{proof}
\runs{tests/calibrate/test\_conformal.py::test\_a\_normal\_observation\_alerts\_at\_most\_alpha\_of\_the\_time}
\runs{tests/calibrate/test\_conformal.py::test\_the\_guarantee\_survives\_a\_distribution\_no\_z\_score\_would}
\runs{tests/calibrate/test\_conformal.py::test\_the\_plus\_one\_is\_the\_guarantee\_and\_not\_rounding}

\begin{corollary}[A budget the history cannot express]\label{cor:resolution}
With $n$ calibration points the smallest attainable $p$-value is $1/(n+1)$. A budget of $b$ false alarms over
$N$ runs asks for $\alpha = b/N$; if $b/N < 1/(n+1)$ no threshold honours it, and \textsc{Prama} refuses the
budget with the arithmetic rather than silently alerting at a higher rate.
\end{corollary}
\runs{tests/calibrate/test\_select.py::test\_a\_business\_budget\_becomes\_a\_statistical\_level}

Exchangeability fails under drift. \textsc{Prama} offers three responses and is explicit about each one's
price: \emph{recency weighting} (a weighted $p$-value with a half-life, which is no longer exact, says so, and
reports its reduced effective sample size); \emph{adaptive levels} in the style of adaptive conformal
inference \cite{gibbs2021}, $\alpha_{t+1} = \alpha_t + \gamma(\alpha - \mathrm{err}_t)$ with $\gamma = 0.02$,
which controls the long-run alert rate rather than each test; and \emph{conditioning} on a declared
seasonal calendar. Adaptation cannot help when every $p$-value sits at its floor, and the code detects that
state rather than adapting forever.
\runs{tests/calibrate/test\_conformal.py::test\_a\_weighted\_p\_value\_does\_not\_claim\_to\_be\_exact}
\runs{tests/calibrate/test\_conformal.py::test\_recency\_weighting\_costs\_resolution\_and\_the\_cost\_is\_reported}
\runs{tests/calibrate/test\_conformal.py::test\_the\_long\_run\_alert\_rate\_converges\_to\_the\_budget\_under\_drift}
\runs{tests/calibrate/test\_conformal.py::test\_adapting\_cannot\_help\_when\_every\_p\_value\_is\_at\_its\_floor}

\begin{table}[t]
\centering\small
\begin{tabular}{@{}lrrrrr@{}}
\toprule
\textbf{Regime} & \textbf{plain} & \textbf{seasonal} & \textbf{weighted} & \textbf{adaptive} & \textbf{changepoint} \\
\midrule
stationary    & 0.0040 & 0.0152 & 0.0096 & 0.0016 & 0.0040 \\
seasonal      & 0.0672 & 0.0232 & 0.0696 & 0.0152 & 0.0672 \\
level shift   & 0.1112 & 0.0904 & 0.0208 & 0.0168 & 0.0784 \\
regime switch & 0.0112 & 0.0520 & 0.0136 & 0.0080 & 0.0112 \\
bursty        & 0.0104 & 0.0192 & 0.0080 & 0.0064 & 0.0104 \\
\bottomrule
\end{tabular}
\caption{Calibration error (mean absolute gap between nominal and realised alert rate over the levels
0.01, 0.02, 0.05, 0.10 and 0.20) by drift regime and mechanism, on synthetic generators at the default 500
observations per regime. Printed by \texttt{tests/monitor/test\_benchmark.py} in our run; the target is 0.02.}
\label{tab:calib}
\end{table}

\Cref{tab:calib} is the result as a grid rather than a number, because no single mechanism is the answer
and one number would hide that. Conditioning on season repairs seasonality and cannot repair a level shift;
forgetting (weighting) nearly repairs a level shift and makes seasonality worse; in this run the adaptive level
meets the target in every regime. What the test asserts is weaker than the table: that every regime has
\emph{some} mechanism meeting the target, that every mechanism is calibrated on the stationary regime, and
that plain conformal fails on seasonal data.
\runs{tests/monitor/test\_benchmark.py::test\_every\_regime\_has\_a\_mechanism\_that\_holds\_the\_promise}
\runs{tests/monitor/test\_benchmark.py::test\_the\_stationary\_case\_is\_calibrated\_by\_every\_mechanism}
\runs{tests/monitor/test\_benchmark.py::test\_plain\_conformal\_fails\_on\_seasonal\_data}

End to end, a monitor declaring a budget of 0.05 realised 0.0467 over 300 judged days in our run, with a
calibration error of 0.0061 on its curve. The test asserts only that more than 200 days were judged and that
the realised rate is within 0.03 of the budget; the two figures are printed, not asserted.
\partly{\texttt{tests/monitor/test\_wave7\_demo.py::test\_the\_alert\_volume\_tracks\_the\_declared\_budget}
asserts a looser bound than the figures quoted.}

\subsection{Selection across many monitors}

With ten thousand monitors each at $\alpha = 0.05$, five hundred false alarms a day is the expected outcome.
\textsc{Prama} selects which $p$-values to raise with Benjamini--Hochberg \cite{bh1995}, or with
Benjamini--Yekutieli \cite{by2001} when dependence is arbitrary, over a hierarchy domain $\to$ dataset $\to$
attribute $\to$ check; a parent's $p$-value is combined from its children by Simes' method, and a wholesale
failure of most checks under one parent is rolled up into a single finding.
\runs{tests/calibrate/test\_select.py::test\_bh\_controls\_the\_false\_discovery\_rate}
\runs{tests/calibrate/test\_select.py::test\_by\_is\_strictly\_more\_conservative\_than\_bh}
\runs{tests/calibrate/test\_select.py::test\_a\_wholesale\_failure\_becomes\_one\_finding\_rather\_than\_four\_hundred}
\stated{a proof that the hierarchical procedure controls FDR at the declared level; the test measures flat BH,
and the roll-up is a presentation rule, not part of the error guarantee.}

\absent{calibration on real financial data. Every figure in this section is from synthetic generators;
the claim that the guarantees hold on a bank's metric history is the central open question of this part of
the work.}

% =============================================================================
\section{Metadata, correlation and fitness for purpose}\label{sec:metadata}

Owners do not only declare structure. They write descriptions, tag fields (``mandatory'', ``allowed
values'', ``PII''), and map attributes to glossary terms. Two further derivations use that material.

\begin{definition}[Metadata-implied rule]\label{def:metarule}
A metadata template field of a given kind (flag, choice, list, number, \ldots) may carry a rule pattern with
placeholders. Setting the field's value on an attribute instantiates the pattern with the value as a
\emph{typed literal} and submits the result as a proposal with origin metadata. Changing the value retracts
the earlier proposal and keeps its history.
\end{definition}

\begin{proposition}[A value cannot become syntax]\label{prop:injection}
Instantiation quotes each value as a literal of the placeholder's type, and a pattern may use only
placeholders its field kind can fill. A value such as \texttt{x') OR (1=1} is therefore a string literal in
the resulting control, not a change to its predicate.
\end{proposition}
\runs{tests/semantic/test\_metadata.py::test\_metadata\_proposes\_its\_rules\_and\_retracts\_them\_when\_changed}
\runs{tests/semantic/test\_metadata.py::test\_a\_value\_cannot\_inject\_pql}
\runs{tests/semantic/test\_metadata.py::test\_a\_rule\_may\_only\_use\_placeholders\_its\_kind\_can\_fill}

\begin{definition}[Correlation]\label{def:correlate}
Attributes are grouped by shared meaning, trying in order a business concept, then a glossary term, then a
semantic type; each attribute joins only its strongest group. Within a group, if exactly one dataset owns
the value as a unique key, every other member is proposed a \pql{REFERENCES} check to it. Generic semantic
types (currency, dates, country, boolean) group for consistency only and never yield a reference.
\end{definition}
\runs{tests/semantic/test\_correlate.py::test\_a\_shared\_term\_groups\_and\_the\_key\_owner\_is\_referenced}
\runs{tests/semantic/test\_correlate.py::test\_no\_owner\_or\_two\_owners\_means\_no\_reference}
\runs{tests/semantic/test\_correlate.py::test\_a\_generic\_type\_groups\_for\_consistency\_but\_never\_proposes\_a\_reference}
\runs{tests/semantic/test\_correlate.py::test\_an\_attribute\_joins\_only\_its\_strongest\_group}

Query history adds a hint that meaning cannot: datasets the warehouse reads in the same query, with a column
in common and no declared relationship. Where that column is a key on one side, the likely \pql{REFERENCES}
is suggested to a steward. A shared name is not a shared meaning, so the hint never becomes a proposal, and
like every usage signal it never enters a score.
\runs{tests/semantic/test\_queried\_together.py::test\_datasets\_read\_together\_suggest\_the\_relationship}

\begin{plain}
\pt If two columns in two datasets are both ``the customer's LEI'', and one dataset is the customer master
where the LEI is unique, then every LEI elsewhere should exist in that master. If nobody, or two datasets,
claim to own it, \textsc{Prama} proposes nothing, because it cannot tell which way the check should point.
\end{plain}

\paragraph{Fitness search.} A question in plain words --- ``trade amounts in USD'', ``who is the customer
and where are they registered, for sanctions screening'' --- is answered by ranking \emph{datasets}, not rows,
over their descriptions, metadata and attribute profiles: with Okapi BM25 \cite{robertson2009}
($k_1 = 1.5$, $b = 0.75$) when no embedding model is configured, and with embeddings when one is. The result
names the matched evidence (which attributes and fields matched) so that the ranking can be checked. A
discover-purpose model may re-rank and explain; it does not decide what exists.
\runs{tests/semantic/test\_fitness.py::test\_without\_an\_embedding\_model\_bm25\_ranks\_and\_names\_its\_evidence}
\runs{tests/semantic/test\_fitness.py::test\_with\_an\_embedding\_model\_profiles\_are\_embedded\_once\_and\_again\_on\_change}
\absent{any retrieval-quality measurement (precision at $k$, nDCG) of fitness search; and correlation from
data-side signals such as value overlap and query co-access.}

\begin{worked}
\wk{governance from metadata}
Case study 7 declares three datasets and has the owners describe them: \texttt{customers.segment} allowed
values \texttt{RETAIL, SME, CORPORATE} and mandatory; \texttt{payments.amount} minimum 0; the two LEI columns
mapped to one glossary term; dataset keys. Nine rules follow as proposals, a reviewer accepts them, and
correlation adds the reference from \texttt{accounts.customer\_lei} to \texttt{customers.lei}, and also
reports that the two columns sharing the term disagree on their PII flag. All four
planted defects are found: six segments \texttt{'SME '} with a trailing space, four accounts with no LEI and
three with an LEI no customer has (together the 7 integrity failures), and five negative payments. Asked
``who is the customer and where are they registered, for sanctions screening'', fitness search returns
\emph{Customers}, naming \texttt{legal\_name}, \texttt{customer\_id} and \texttt{lei} as its evidence.
\end{worked}

% =============================================================================
\section{Evaluation}\label{sec:eval}

We report two kinds of evidence, and one absence. The case studies show the whole pipeline on planted
defects, with both columns --- what was planted and what was found --- printed by the run. The benchmark
measures bounds and ablations on a labelled corpus. Neither compares \textsc{Prama} with another system.

\subsection{Eight case studies}

Each study fabricates seeded banking data, declares an estate in business terms, lets $\Gam$ derive the
controls, runs them, and prints the planted defects against the findings. The rule the studies follow is
that every planted defect is stated before \textsc{Prama} is pointed at anything, in the generator and in the
README, including the ones \textsc{Prama} will not find. We ran all eight
(\texttt{python run.py --no-serve} in each study directory, against a scratch application database);
\Cref{tab:studies} is transcribed from those runs.

\begin{table}[t]
\centering\small
\begin{tabular}{@{}L{0.30\linewidth}rrrrr@{}}
\toprule
\textbf{Study} & \textbf{Planted} & \textbf{Found} & \textbf{Not est.} & \textbf{Not claimed} & \textbf{Controls} \\
\midrule
1 Trading book (SQLite)            & 9 & 4 & 2 & 3 & 65 \\
2 Feeds: CSV, Parquet, JSON Lines  & 9 & 7 & 1 & 1 & 58 \\
3 Mixed estate, relationships      & 4 & 4 & 0 & 0 & 59 \\
4 Expressions and plugins          & 5 & 4 & 1 & 0 & 27 \\
5 DQ delegates                     & 4 & 4 & 0 & 0 & 14 \\
6 Month-end close                  & 4 & 3 & 0 & 0 & 18 \\
7 Governance from metadata         & 4 & 4 & 0 & 0 & 10 \\
8 From code to impact              & 2 & 2 & 0 & 0 & 7 \\
\midrule
Total                              & 41 & 32 & 4 & 4 & \\
\bottomrule
\end{tabular}
\caption{Planted defects against outcomes, from our runs of the eight case studies. \emph{Not est.}: reported
as not established (a screen, or, when these runs were made, a freshness control on a plain table). \emph{Not claimed}: stated in the study
as outside what the declared controls can see, with the reason. In study 6 the fourth planted item is a
timing difference that is correctly \emph{not} a break. \emph{Controls} is the number executed on the main
source; study 3 runs a further 22 on a second source.}
\label{tab:studies}
\end{table}

\paragraph{What the table shows.} Every defect a declared control can see was found, on the rows planted,
with counts that match: 60 missing notionals, 7 negative quantities, 400 duplicated rows from a resent file,
37 orphan instruments, 18 settlement dates before trade dates, 31 half-cent rounding disagreements. The four
``not established'' outcomes are the screens of \Cref{sec:screen} (ISIN, LEI and a bank's book-code check
character) and the freshness control of \Cref{sec:gamma} as it then was; in none of them did \textsc{Prama}
report a pass. Rerun now, study 1 generates no freshness control and says why, since its tables declare no
arrival column.
The four ``not claimed'' are a relationship not declared in study 1 (two defects) and a relation between two
dates nobody declared, and in study 2 a truncated file whose trailer check runs at arrival and is not yet
wired into the control runner.

\paragraph{What the table does not show.} The defects were planted by the author of the system that finds
them, in data generated for the purpose, so the table demonstrates mechanism, not detection quality. Its
most useful columns are the last two: a tool with no false negatives on data it was built against tells a
reader nothing, while a tool that names what it did not catch, and why, tells the reader what it is. The
per-study counts are printed by each \texttt{run.py} and stated in each README; the test suite asserts only
that each study runs to completion.
\partly{\texttt{tests/casestudies/test\_they\_still\_run.py::test\_it\_runs\_to\_completion} asserts exit status
0 for each study; the planted-versus-found counts are not asserted.}

\paragraph{What building them found.} Running the studies against data somebody had to look at found six
defects in \textsc{Prama}, all of the kind this paper is about --- artefacts that build, validate and look right
while being wrong. $\Gam$ generated \pql{MATCHES /None/} for every free-text column, because a value domain
read back from JSON kept its kind as a string and the ``is this constrained?'' test passed for everything;
controls were lowered without their code lists; a regular expression would not bind against a date column on
a strict engine; SQLite was believed unable to do regular expressions; \pql{ROUND} double-rounded on DuckDB;
and an unknown function compiled straight to SQL while the reference interpreter returned unknown. Study 8
found two more, both fixed in the most recent revision: a lineage reconciliation proposed across a filter,
which reported 118 cancelled trades as breaks, and a proposed reconciliation with no tolerance, which the
engine refused to run.

\subsection{The labelled defect benchmark}

\texttt{prama bench run --seed 42} builds a corpus of 28 scenarios of 200 rows, one per defect class, across
six families (structural, content, statistical, relational, temporal, semantic) and four difficulties
(obvious 7, ordinary 9, subtle 11, adversarial 1), and labels only the rows that actually changed. A detection
counts only when dataset, column and window all match. Our run:

\begin{center}\small
\begin{tabular}{@{}llrrrr@{}}
\toprule
\textbf{Detector} & \textbf{Kind} & \textbf{Found} & \textbf{Precision} & \textbf{Recall} & \textbf{F1} \\
\midrule
detect-nothing      & bound    & 0/28  & --   & 0.00 & -- \\
alert-on-everything & bound    & 28/28 & 0.08 & 1.00 & 0.14 \\
schema-only         & ablation & 2/28  & 0.67 & 0.07 & 0.13 \\
patterns-only       & ablation & 5/28  & 0.83 & 0.18 & 0.29 \\
statistics-only     & ablation & 7/28  & 0.88 & 0.25 & 0.39 \\
\midrule
prama-declared      & system   & 10/28 & 0.77 & 0.36 & 0.49 \\
\bottomrule
\end{tabular}
\end{center}

The blind spots matter more than the F1: statistics-only found nothing in the relational family, and
patterns-only nothing in the statistical, relational, temporal or semantic families. A detector's blind
spots are the argument for whatever covers them, and they do not appear in an aggregate score.
\runs{tests/bench/test\_corpus.py::TestReproducibility::test\_the\_same\_seed\_gives\_the\_same\_corpus}
\runs{tests/bench/test\_corpus.py::TestTheLabelsFollowWhatHappened::test\_a\_no\_op\_injection\_plants\_no\_label}
\runs{tests/bench/test\_baselines.py::TestTheBoundsAreTheAxes::test\_detecting\_nothing\_has\_no\_precision\_rather\_than\_zero}
\runs{tests/bench/test\_baselines.py::TestTheAblationsAnswerWhichClaimDoesTheWork::test\_no\_ablation\_comes\_near\_perfect\_recall}
\runs{tests/bench/test\_scoring.py::TestTheMatchIsExact::test\_the\_wrong\_column\_does\_not}
\runs{tests/cli/test\_bench\_cli.py::TestRun::test\_it\_names\_the\_baselines\_that\_were\_not\_run}

\paragraph{The system.} The last row is \textsc{Prama}'s declared path: controls $\Gam$ derives from an owner's
declaration of the dataset, run on each window by the reference interpreter. The declaration is written from the
schema's documented domain, not from the defect classes, and is not revised to catch what it misses. It uses
no relationships (the corpus has one dataset), no history and no monitors, so it is one detector of several.
It finds every structural defect, four of seven content defects and one statistical one. It finds nothing
relational, temporal or semantic, which are the families that need a second dataset, a clock, or a meaning
the schema does not state. Its fairness test, that no clean window raises an alert, found a defect in the
corpus itself: the generator's IBANs had made-up check digits, so a correct IBAN validator alerted on every
window. The generator now computes them.
\runs{tests/bench/test\_declared.py::test\_clean\_windows\_raise\_nothing}
\runs{tests/bench/test\_declared.py::test\_the\_declared\_path\_scores\_and\_names\_what\_it\_cannot\_see}

\paragraph{The absence.} None of the fifteen external systems the benchmark names (Great Expectations, Soda
Core, Deequ, dbt tests, HoloClean/Raha/Baran/ZeroED and others) was run; the command prints them as not run.
We regard publishing bounds, ablations and our own score without a comparison as better than publishing a
comparison configured by the party that benefits from it. It means this paper claims that the declared path
beats every single-technique ablation here, and makes \emph{no} claim relative to any other product.
\absent{any external baseline; the banking-scale corpus; a blind evaluation with stewards (the blinding mechanism exists: \texttt{tests/bench/test\_shadow.py::TestTheBlinding::test\_what\_a\_steward\_sees\_carries\_no\_system}).}

\subsection{Measurements of the system itself}

The following were measured in our runs and concern \textsc{Prama}'s own overhead, not detection. The test
packages cited in this paper (1,693 tests across \texttt{backend}, \texttt{evidence}, \texttt{recon},
\texttt{derive}, \texttt{lineage}, \texttt{architecture}, \texttt{calibrate}, \texttt{score}, \texttt{llm},
\texttt{induce}, \texttt{propose}, \texttt{semantic}, \texttt{codeintake}, \texttt{delegates} and \texttt{bench})
passed, with one strict expected failure (the freshness control of \Cref{sec:gamma}, since closed), in 25
seconds. Fusion ran
400 generated controls in 160 scans. The backend conformance suite passed on DuckDB, SQLite, the reference
interpreter and PostgreSQL 16.

% =============================================================================
\section{The claims register}\label{sec:register}

A paper that argues that beliefs about data must be justified owes the same to its own claims. The
register below lists every formal claim with its state: \textbf{runs} (a named test exercises it),
\textbf{in part} (something weaker runs, as stated at the claim), \textbf{stated} (mathematics or a
specification that constrains the system and that the system does not execute), and \textbf{not built}.

\begin{center}
\small
\begin{longtable}{@{}L{0.42\linewidth}L{0.12\linewidth}L{0.38\linewidth}@{}}
\toprule
\textbf{Claim} & \textbf{State} & \textbf{Where, or why not} \\
\midrule
\endhead
\multicolumn{3}{@{}l}{\emph{Declarations and $\Gam$}} \\
Relationships, thirteen kinds, never auto-confirmed (\Cref{def:rel}) & runs & \texttt{test\_services.py} \\
Determinism and identity stability (\Cref{prop:det}) & runs & \texttt{test\_generator.py} \\
Provenance (\Cref{prop:prov}) & runs & \texttt{test\_generator.py} \\
Round trip (\Cref{prop:roundtrip}) & runs & \texttt{test\_generator.py} \\
Coalescing (\Cref{prop:coalesce}) & runs & \texttt{test\_generator.py} \\
Refusal, not omission (\Cref{prop:refusal}) & runs & \texttt{test\_relationships.py} \\
Honest and flattering coverage (\Cref{prop:flatter}) & runs & \texttt{test\_coverage.py} \\
Accuracy needs a relationship (\Cref{prop:accuracy}) & runs & \texttt{test\_wave6\_acceptance.py}; 82/71\% figures prose \\
Only a declaration may skip review (\Cref{prop:review}) & runs & \texttt{test\_queue.py}, \texttt{test\_provenance.py} \\
A person activates every control (\Cref{sec:proposals}) & in part & one activation path, by inspection; untested \\
Corroboration merges (\Cref{prop:corrob}) & runs & \texttt{test\_queue.py}, \texttt{test\_utility.py} \\
A rhythm yields a freshness verdict & runs & \texttt{test\_freshness.py}, \texttt{test\_generator.py} \\
PQL parity and roll-forward forms; hierarchy cycles & not built & comparison specifications only \\
\multicolumn{3}{@{}l}{\emph{PQL, plans and verdicts}} \\
Plan identity closes over code lists and validators (\Cref{prop:closure}) & in part & exclusions untested \\
Nothing scanned is not a pass (\Cref{prop:empty}) & runs & \texttt{test\_two\_stage.py} \\
Unknown default is conservative (\Cref{prop:unknown}) & runs & \texttt{test\_parser.py}, conformance \\
Screen soundness (\Cref{prop:screen}) & runs & \texttt{test\_two\_stage.py}, \texttt{test\_control\_run.py} \\
Residual runs on the pushdown path & in part & reference interpreter only \\
Cross-engine agreement (\Cref{prop:conform}) & runs & 25 cases; 250 generated; PostgreSQL when configured \\
Semantic equivalence of compiler and interpreter & stated & tested, not proved \\
Refusal is conformance (\Cref{prop:refuse}) & in part & RE2 refusal untested \\
Fusion preserves answers and saves scans (\Cref{prop:fusion}) & runs & \texttt{test\_fuse.py}, \texttt{test\_throughput.py} \\
Custom SQL is one read-only query & runs & \texttt{test\_custom\_sql.py} \\
A delegate cannot decide (\Cref{prop:delegate}) & runs & \texttt{test\_delegates.py} \\
Delegate sandbox isolation & runs & \texttt{test\_sandbox.py}; namespace where the host allows \\
Exact arithmetic & in part & operators exact; thresholds float \\
\multicolumn{3}{@{}l}{\emph{Evidence}} \\
Tamper evidence (\Cref{thm:tamper}) & runs & \texttt{test\_ledger.py}, \texttt{test\_independent\_verifier.py} \\
External anchoring of chain heads & runs & \texttt{test\_anchor.py}; RFC 3161, verified by OpenSSL \\
Merkle promotion (\Cref{def:merkle}) & runs & \texttt{test\_ledger.py} \\
Independent verification (\Cref{prop:independent}) & runs & \texttt{test\_independent\_verifier.py} \\
Erasure keeps the chain (tombstones) & runs & \texttt{test\_independent\_verifier.py} \\
Mixed format versions verify & runs & \texttt{test\_format\_versions.py} \\
Records under 2 KB & runs & \texttt{test\_ledger.py}; 848 bytes is prose \\
Replay classifies (\Cref{prop:replay}) & runs & \texttt{test\_replay.py}, \texttt{test\_end\_to\_end.py} \\
\multicolumn{3}{@{}l}{\emph{Reconciliation}} \\
Tolerance laws (\Cref{prop:tolerance}) & runs & \texttt{test\_reconciliation.py} \\
Exact match first (\Cref{lem:exactfirst}) & runs & \texttt{test\_reconciliation.py} \\
Price of ignoring timing (\Cref{prop:timing}) & in part & case study 6 run; not asserted \\
Break classes (\Cref{def:breaks}) & runs & \texttt{test\_reconciliation.py} \\
A filter on each side of \pql{RECONCILE} & runs & \texttt{test\_reconcile\_pql.py}, \texttt{test\_joins.py} \\
\pql{CLASSIFY} on \pql{RECONCILE}; signed certificate & not built & \\
\multicolumn{3}{@{}l}{\emph{Lineage}} \\
Parsed, never executed; fallback recorded & runs & \texttt{test\_pyspark.py}, \texttt{test\_sql.py}, \texttt{test\_intake.py} \\
Gold-fixture precision and recall & runs & \texttt{test\_gold\_fixture.py} \\
Strongest path (\Cref{prop:strongest}) & runs & \texttt{test\_graph.py} \\
Attenuation constants calibrated & stated & declared constants, not estimates \\
What a copy justifies (\Cref{prop:copy}) & runs & \texttt{test\_lineage\_controls.py}, \texttt{test\_reconcile\_pql.py} \\
Join predicates invisible (\Cref{prop:join}) & in part & case study 8 run; the repair is asserted \\
Population-dependence edges & in part & SQL and PySpark readers; a pandas \texttt{merge} is a gap \\
A change that removes a control's basis fails review & runs & \texttt{test\_review.py} \\
Every driving row finds its match (join checks) & runs & \texttt{test\_reconcile\_pql.py} \\
\multicolumn{3}{@{}l}{\emph{Models}} \\
Architectural guard (\Cref{prop:guard}) & runs & \texttt{test\_layering.py}, with counterfactual \\
Model-free verdict path, as information flow & stated & guard is syntactic \\
No vacuous control reaches a reviewer (\Cref{prop:vacuous}) & runs & \texttt{test\_validate.py}, \texttt{test\_inducer.py} \\
Hash-chained call ledger (\Cref{def:callledger}) & in part & chain tested; CLI untested \\
Evaluation gates activation (\Cref{prop:evalgate}) & runs & \texttt{test\_governance.py} \\
Residency, redaction, budgets & runs & \texttt{test\_providers.py}, \texttt{test\_gateway.py} \\
Quality of model-authored rules & not built & not measured \\
\multicolumn{3}{@{}l}{\emph{Trust and calibration}} \\
Trust choices disagree (\Cref{def:trust}) & runs & \texttt{test\_trust.py} \\
Semiring laws of the trust choices & stated & only one is a semiring \\
Conformal validity (\Cref{thm:conformal}) & runs & \texttt{test\_conformal.py} \\
Unexpressible budgets refused (\Cref{cor:resolution}) & runs & \texttt{test\_select.py} \\
Drift mechanisms and their costs & runs & \texttt{test\_conformal.py} \\
Some mechanism calibrated per regime (\Cref{tab:calib}) & runs & \texttt{test\_benchmark.py} \\
Declared budget realised end to end & in part & looser bound asserted than quoted \\
BH and BY control FDR & runs & \texttt{test\_select.py} \\
Hierarchical FDR guarantee & stated & roll-up is presentation \\
Calibration on real financial data & not built & synthetic only \\
\multicolumn{3}{@{}l}{\emph{Metadata and search}} \\
A value cannot become syntax (\Cref{prop:injection}) & runs & \texttt{test\_metadata.py} \\
Correlation and key ownership (\Cref{def:correlate}) & runs & \texttt{test\_correlate.py} \\
Fitness search names its evidence & runs & \texttt{test\_fitness.py} \\
Queried-together hints stay hints & runs & \texttt{test\_queried\_together.py} \\
Retrieval quality of fitness search & not built & not measured \\
\multicolumn{3}{@{}l}{\emph{Evaluation}} \\
Benchmark bounds and ablations & runs & \texttt{tests/bench} \\
\textsc{Prama} scored on its own benchmark & runs & \texttt{test\_declared.py}; declared path only \\
External baselines & not built & named, not run \\
Case studies: planted against found (\Cref{tab:studies}) & in part & exit status asserted; counts printed \\
\bottomrule
\end{longtable}
\end{center}

Of seventy-five rows, fifty-three run, eleven run in part, five are stated without being executed, and six
are not built.

The ratio is not flattering and is not meant to be. The claims that run are concentrated where the thesis
lives: derivation, verdict semantics and evidence. The claims that do not are concentrated where a claim is
about the world rather than about the code --- detection quality against other systems, calibration on real
data, whether reviewers accept what models propose --- and those need a bank, a steward and time, none of
which a test suite can supply.

% =============================================================================
\section{Related work}\label{sec:related}

\paragraph{Rule frameworks and validation.} Deequ \cite{deequ} and TFX data validation \cite{tfdv} made
declarative assertions over large datasets practical, and open frameworks in the same family (Great
Expectations, Soda, dbt tests) are widely deployed. They take the assertion as the primitive. \textsc{Prama}
takes the \emph{declaration} as the primitive and derives assertions, which is what makes coverage measurable
and a control's existence explainable. Data contracts \cite{odcs} move in the same direction, though
contract tooling frequently records obligations without enforcing them \cite{contractenforce}; \textsc{Prama}
imports ODCS quality blocks where they state a rule it has, and refuses the rest by name. SHACL
\cite{shacldq,shaclreview} is the closest standard to a declaration-first model for graph data.

\paragraph{Discovery and cleaning.} Profiling and dependency discovery \cite{metanome,fdeval,dcdiscovery,
conformance} find constraints the data satisfies; error detection and repair \cite{holoclean,rahabaran,
zeroed} find cells that are probably wrong. Both are sources of \emph{proposals} in \textsc{Prama}, ranked
below a declaration because a column that happens to be unique today is not a declared key. Example-driven
induction follows the label-efficient line of Raha and Baran \cite{rahabaran} and weak supervision
\cite{snorkel}.

\paragraph{Models as rule authors.} Recent work uses language models to generate data quality rules
\cite{llmdqr,llmtabular,tuwien}. Our contribution here is a boundary rather than a generator: the
counterfactual gate (\Cref{prop:vacuous}), the evaluation gate on activation (\Cref{prop:evalgate}), and a
structural argument that a model can influence a verdict only through a person's activation.

\paragraph{Calibration and multiple testing.} Conformal prediction \cite{vovk2005} gives distribution-free
false-alarm guarantees under exchangeability; adaptive conformal inference \cite{gibbs2021} trades per-test
validity for long-run rate control under drift; conformal anomaly detection with FDR control is an active
area \cite{fdrnovelty,nonconform}. Benjamini--Hochberg \cite{bh1995} and Benjamini--Yekutieli \cite{by2001}
control the false discovery rate across tests. We apply these off the shelf; the contribution is honest
degradation (refusing a budget the history cannot express) and reporting a grid rather than a number.

\paragraph{Provenance and lineage.} Provenance semirings \cite{green2007} give the algebra of which our trust
choices are instances, or near-instances (\Cref{sec:trust}). OpenLineage \cite{openlineage} standardises the
exchange of lineage events; \textsc{Prama} emits it and reads SQL, Python and BI models itself, and our
negative result (\Cref{prop:join}) concerns any lineage model that records value flow only.

\paragraph{Tamper-evident records.} Hash-linked timestamping \cite{haber1991}, Merkle trees \cite{merkle1987},
tamper-evident logging \cite{crosby2009} and certificate transparency \cite{laurie2013} are the foundations of
the evidence ledger. The ledger's contribution is modest and practical: a verifier that shares no code with the
writer, held equal to it by a test, and an RFC 3161 anchor after every run that the auditor verifies with
their own tools.

\paragraph{Regulation and practice.} BCBS 239 \cite{bcbs239} frames the obligation; fines for reporting
errors \cite{ateam2024} price it; market surveys \cite{gartner2025adq,datakitchen2026,dqframeworks} describe the
fragmented tooling that the thesis responds to.

% =============================================================================
\section{Limitations and how this could be wrong}\label{sec:limits}

\paragraph{Do owners declare?} The thesis rests on business owners writing declarations. The case studies
declare on the owners' behalf. Nothing here measures whether a real steward, in a real week, produces the
declarations that make $\Gam$ useful. If they do not, derivation has nothing to derive from, and the system
falls back to mining and induction --- the proposals this paper ranks lowest.

\paragraph{Planted defects are not defects.} Every defect in the evaluation was planted by the author, and the
benchmark scores only \textsc{Prama}'s declared path against bounds and ablations. The paper therefore makes
no claim of superior detection over any other product. What it
claims is narrower: that a verdict \textsc{Prama} reports is justified in the sense of \Cref{sec:thesis}, and that
what it cannot justify it reports as not established.

\paragraph{Synthetic calibration.} The conformal guarantees are theorems under exchangeability, and the drift
mechanisms are measured on generators the author wrote. A bank's metric history may violate the assumptions in
ways none of the five regimes represents.

\paragraph{Scale.} The declaration model and $\Gam$ have not been exercised on an estate of thousands of
datasets; the throughput figures measure \textsc{Prama}'s own overhead against an in-memory engine, not a
warehouse.

\paragraph{Connectors and engines.} Conformance runs on DuckDB, SQLite and PostgreSQL. Other dialects reached
through JDBC are code-complete but not verified against live services, and a Snowflake connector has never met
an account; the repository says so wherever it describes them.

\paragraph{The guard is syntactic.} \Cref{prop:guard} is a necessary condition checked by co-occurrence of
names. It would not detect a model's output passed through neutral names across modules. The structural
argument --- a verdict is a function of an activated plan and the data --- is the stronger defence, and it is
argued, not checked.

\paragraph{Evidence trusts its anchors.} With anchoring off, a party holding the HMAC key can rewrite a chain
and re-sign its head. With it on, the records written since the last anchor remain rewritable until the next;
and the anchor is only as independent as the time-stamp authority chosen.

\paragraph{What would falsify the thesis.} A declared estate on which derived controls miss defects that
hand-written controls catch, at comparable effort; a deployment in which $\IND$ is so frequent that operators
treat it as a pass; or a replay that reports \textsf{identical} for a run whose data changed. Each is a
measurable outcome, and the first two need real users.

% =============================================================================
\section{Conclusion}\label{sec:conclusion}

A data quality claim is a belief, and a belief is worth acting on when it is justified. We turned that into
three constraints --- controls derived from declarations, verdicts from a deterministic engine, evidence that
verifies without its author --- and showed what each requires: a total, provenance-carrying map from
declarations to controls with a coverage measure that does not flatter; a verdict set in which ``the SQL found
nothing'' can be reported as what it is; and a record an auditor can check with the standard library. On that
base, reconciliation acquires a tolerance semantics and a timing law, lineage acquires a precise statement of
which controls a copy justifies, and models acquire a boundary they cannot cross except through a person.

The register counts fifty-three claims that run and twenty-two that do not, or do so only in part. The ones
that run are the ones the thesis depends on. The ones that do not are, for the most part, claims about people
and data that the author does not own, and they are where the work goes next.

\bibliographystyle{plainurl}
\bibliography{references}

\end{document}
